\documentclass[../main.tex]{subfiles}
\begin{document}
%==========================================TIÊU ĐỀ TẬP========================================
\begin{table}[h]
  \begin{adjustwidth}{-1cm}{}
    \begin{tabular}{>{\centering\arraybackslash}p{6cm}|>{\raggedright\arraybackslash}p{12.5cm}}
      \multirow{5}{6cm}
      % Cột bên trái: ÉPISODE và số tập
      {\\[-40pt]
        \flaregothic\fontsize{20pt}{20pt}\selectfont \centering
        {\contour{errorfive}{\textcolor{black}{ÉPISODE}}}\color{black}\\
        \burbank\fontsize{72pt}{72pt}\selectfont
        \includegraphics[height=3.12359cm]{../../../../Image/Book?/number/num33.png}
        \\[18pt]
        % \flaregothic\fontsize{14pt}{14pt}\selectfont
        % \color{epattr}BIENVENUE, \uline{\color{Banpire} Mel} !
        % \flaregothic\fontsize{18pt}{18pt}\selectfont
        % \color{epattr}ÉPISODE PERDU
        \color{black}
      }
       &
      % Dòng thứ nhất: tên tập tiếng Nhật
      {\fotskip\fontsize{18pt}{18pt}\selectfont \ruby{解}{と}かれた\ruby{記}{き}\ruby{憶}{おく}}
      \\[6pt]
      % Dòng thứ hai: tên tập tiếng Anh
       & {\cafeteria\fontsize{18pt}{18pt}\selectfont The Last Memories}                                                                                                                                     \\[4pt]
      % Dòng thứ ba: tên tập tiếng Việt
       & {\vnmsans\fontsize{18pt}{18pt}\selectfont Mảnh ký ức cuối cùng}                                                                                                                                    \\[8pt]
      % Dòng thứ tư: nhân vật xuất hiện
       & {\caviar{\fontsize{14pt}{14pt}\selectfont Nhân vật: \emp{kuon1}, \emp{kaigou1}, \emp{suzuki1}, \emp{sora}, \emp{miko}, \emp{azki},\emp{suisei}, \emp{shion}, \emp{ayame}, \emp{choco}, \emp{marine}}}
      % \\[6pt]
      % Dòng cuối cùng: Thời gian diễn ra của tập (thường sẽ là ngày phát hành)
      % & {\continuum\fontsize{16pt}{16pt}\selectfont Ngày phát hành: 11/06/2022}\\[8pt]
    \end{tabular}
  \end{adjustwidth}
\end{table}
%===========================================NỘI DUNG==========================================
\color{black}

\pov{Hikawa Kuon}

Ngày lễ Valentine đã trôi qua được vài ngày, nhưng dư vị ngọt ngào của sô-cô-la và những tiếng cười náo nhiệt dường như không thể xua tan đi bầu không khí đặc quánh đang bao trùm lấy thị trấn Aogami. Ngược lại, tôi cảm giác như có một thứ gì đó đang dần cô đặc lại, lẩn khuất sau những bóng râm của hàng thông và len lỏi vào từng hơi thở của màn đêm.

Những hiện tượng lạ đang gia tăng với tần suất chóng mặt.

Không còn là những thoáng qua mờ ảo nữa. Bóng dáng của những đứa trẻ bắt đầu xuất hiện thường xuyên hơn ở những góc hành lang tối, hay thoáng hiện sau những gốc cây bàng già cỗi trong sân trường. Chúng không làm gì cả, chỉ đứng đó, lặng lẽ nhìn theo chúng tôi bằng những đôi mắt sâu thẳm, trống rỗng nhưng chứa chan một nỗi mong mỏi vô hình.

Suzuki là người nhạy cảm nhất. Con bé liên tục nói với tôi về việc mình luôn cảm thấy có ``những ánh nhìn chằm chằm'' từ phía sau lưng, ngay cả khi đang ở trong phòng học. Mỗi khi quay lại, con bé chỉ thấy những hạt bụi li ti nhảy múa dưới ánh đèn huỳnh quang, nhưng cảm giác lạnh sống lưng thì vẫn vẹn nguyên.

Ngay cả Mari và Nanase cũng bắt đầu có những cơn déjà vu mạnh đến mức đáng lo ngại. Mari-san đôi khi đột ngột dừng bút giữa chừng, nhìn trân trân vào bức vẽ khu rừng như thể cô ấy vừa nhớ ra một chi tiết nào đó chưa từng được vẽ xuống. Còn Nanase-san, cô ấy hay thẫn thờ nhìn ra phía cầu thang khu nhà C, nơi mà những linh hồn trẻ thơ dường như thường xuyên tụ tập.

Tôi tựa lưng vào cửa sổ, mắt hướng về phía khu rừng Aogami đang thẫm màu dưới ánh hoàng hôn tà. Một cảm giác bất an dâng lên trong lòng.

Những linh hồn đó\dots\ chúng không còn đơn thuần là những kẻ quan sát hay những bóng ma của quá khứ nữa. Chúng đang chờ đợi. Một sự chờ đợi mòn mỏi kéo dài hơn mười ba thế kỷ, giờ đây đã chạm đến điểm giới hạn. Chúng đang chờ một hành động, một lời công nhận, hay có lẽ là một ``chỗ dừng chân'' cuối cùng để có thể buông bỏ tất cả.

``\textit{Chúng không còn muốn trốn tránh nữa\dots\ chúng đang đợi chúng ta tìm thấy chúng,}'' tôi lầm bầm, cảm nhận được hơi lạnh của mặt kính cửa sổ thấm vào da thịt.

Thời gian của sự im lặng đã kết thúc. Vòng lặp này cần một dấu chấm hết thực sự, trước khi những mảnh ký ức cuối cùng tan biến hoàn toàn vào hư không.

Trong lúc tôi còn đang mải suy nghĩ, một giọng nói lanh lảnh bất chợt vang lên:

``\textit{Kuon, kẻ bị vận mệnh chọn lọc, ngươi có cảm thấy dòng chảy của thời không đang gào thét trong lồng ngực không?}''

Tôi quay lại thì thấy Shion đứng chắn ngay lối đi, hai tay khoác lên chiếc áo khoác đen dài tung bay như cánh dơi, đôi mắt cậu ẩn sau lớp kính áp tròng đỏ rực. Mái tóc bạch kim được vuốt ngược, để lộ chiếc trâm cài hình lưỡi hái tử thần lấp lánh.

``Shion-san? Cậu làm gì ở đây?''

``Ta đang thực thi nhiệm vụ cao cả mà Hội Bí Ẩn giao phó,'' cô nàng nghiêng đầu, giọng khàn đặc quánh drama. ``Trong bóng tối của hành lang trường học, ta đã thấy \textit{tiếng gọi} phát ra từ phòng CLB Huyền bí. Một thế lực cổ xưa vừa thức giấc, và nó chỉ đáp ứng với năm linh hồn được chọn.''

Tôi thở dài, xoa trán. ``Cậu lại xem anime gì rồi?''

``Không phải anime!'' Shion thốt lên, giọng cao như kính vỡ. ``Là \textit{thực tại song song}. Ta vừa phát hiện ra cuốn grimoire thất lạc, và nếu ngươi không tới ngay, cả thị trấn Aogami sẽ bị nuốt chửng bởi vòng lặp ký ức.'' Cô ấy giơ tay, đưa ra một mảnh giấy có huy hiệu lúc nhúc mực đỏ. ``Lệnh triệu tập. Ngươi phải tới trước hoàng hôn, nếu không\dots\ chúng ta mất đi cơ hội duy nhất để \textit{khóa} lại cánh cổng.''

Tôi liếc đồng hồhồ, lúc này đã 4:37 chiều. Hoàng hôn ở Aogami thường xuống sớm. Một cảm giác bất an thực sự lần theo xương sống tôi. Có thể đây chỉ là trò chơi tưởng tượng của Shion, nhưng nếu cô ấy đã tìm được thứ gì đó liên quan tới rừng Aogami\dots\ thì có lẽ đây không phải là điều tôi có thể làm ngơ.

``Được rồi,'' tôi gật đầu. ``Dẫn đường.''

Shion nở nụ cười đắc thắng, xoay người, áo choàng quét một vòng đen kịt. ``Theo sát ta, kẻ bị nguyền. Con đường này không dành cho kẻ yếu tim.''

\ct

Chiều hôm đó, tôi tới phòng CLB Huyền bí. Không chỉ có bộ ba chúng tôi, mà Azuki, Shino và Sakura cũng bị lôi theo. Sự ồn ào của cô nàng Phù thủy tự xưng đã thu hút luôn sự tò mò của Nanase, Ayame và Mari đang đi ngang qua hành lang. Thậm chí, tôi còn thoáng thấy bóng dáng Yuko-sensei đứng khoanh tay dựa vào khung cửa, nửa như đang giám sát, nửa như đang nghe ngóng.

Shion đứng chống nạnh trước bàn giáo viên đã bị trưng dụng làm bàn họp, vẻ mặt đắc thắng như vừa tìm ra kho báu của thế kỷ. Trên bàn là một tập tài liệu ố vàng, bìa ngoài được bọc bằng một lớp giấy da cũ kỹ, bên trên là một phong ấn mờ nhạt đã bị xé rách.

``Khá khen cho sự có mặt đông đủ của các ngươi!'' Shion vỗ tay xuống bàn cái *BỐP*, khiến Mari giật bắn mình. ``Nhờ vào thính giác siêu phàm và sự nhạy bén với ma lực của đại pháp sư Shion ta đây, ta đã đánh hơi được thứ này trong cái đống hồ sơ lưu trữ cũ mèm ở góc phòng! Các ngươi biết nó là gì không?''

``Một cuốn sách phép thuật cổ xưa để gọi ra quái vật chăng?'' Ayame nhai kẹo cao su, nhếch mép trêu chọc.

``Sai bét!'' Shion trừng mắt. ``Đây là thứ sẽ giải đáp toàn bộ những gì chúng ta đã thấy ở rừng Aogami!'' Nàng tóc xám-trắng đặt tập tài liệu lên, trông có vẻ như đã bị nhàu nát và ố vàng bởi thời gian. Tiêu đề của nó là: ``Bản ghi chép nghi thức trấn yểm sạt lở thời Aogami - thế kỷ VIII.''

Không khí trong phòng chùng xuống. Cái tên của tập tài liệu gợi lại những ký ức lạnh lẽo của chuyến đi vào rừng lần trước. Shino khẽ nhíu mày, tiến lên một bước: ``Cậu đã đọc nội dung bên trong chưa, Shion-san?''

``Tất nhiên là rồi. Và nó xác nhận chính xác những gì Kuon đã suy luận,'' Shion giở nhẹ từng trang giấy giòn rụm, giọng cô trầm xuống, không còn vẻ cợt nhả. ``Những đứa trẻ đó không phải chết vì tai nạn hay dịch bệnh. Chúng bị hiến tế sống.''

``Hiến tế sống?'' Suzuki thốt lên, tay siết chặt lấy gấu áo tôi. ``Giống như những gì chúng ta đã thấy trong\dots\ trong vòng lặp đó sao?''

``Chính xác là chúng,'' Sakura gật đầu, đôi mắt màu mận chín lướt nhanh qua các dòng chữ cổ. ``Đây này, trong này có ghi:

\txtbx{Đứa trẻ mang dòng máu thanh khiết sẽ được dâng lên để xoa dịu cơn giận của thần núi. Máu thịt hòa vào đất, linh hồn hóa thành sứ giả.}

Bọn họ tin rằng việc chôn sống những đứa trẻ cùng với bùa ngải và lúa rang sẽ khiến chúng bị trói buộc, trở thành vật thế mạng vĩnh viễn cho ngôi làng.''

``Thật tàn nhẫn\dots'' Mari ôm lấy cuốn sổ vẽ của mình, giọng run rẩy. Nanase bên cạnh cũng cau mày, ánh mắt hiện lên vẻ phẫn nộ: ``Nếu những gì họ làm chỉ để bảo vệ ngôi làng một cách mù quáng\dots\ thật không thể dung thứ!''

``Nhưng điều quan trọng nhất không nằm ở đó,'' tôi chỉ tay vào đoạn văn bản ở trang cuối cùng, nơi những nét mực đen nhánh dường như được viết bằng một sự vội vã, hoảng loạn. ``Vấn đề là, tại sao linh hồn chúng vẫn còn ở đây, hiện diện ở khắp thị trấn này?''

``Bởi vì nghi lễ chưa bao giờ được hoàn tất,'' giọng nói của Yuko-sensei vang lên từ phía cửa. Cô bước vào, ánh mắt sắc sảo nhìn lướt qua tập tài liệu. ``Người lớn thời đó đã lừa gạt lũ trẻ, gieo vào đầu chúng niềm tin rằng sự hy sinh của chúng là cao cả. Nhưng sâu thẳm bên trong, họ biết mình đang phạm tội ác. Sự chối bỏ đó khiến các linh hồn không thể siêu thoát.''

Shion gật gù, dùng bút gõ gõ vào trang giấy: ``Đúng vậy! Tài liệu này chỉ ra một cách hóa giải duy nhất. Không phải là thực hiện một buổi lễ cầu siêu bình thường, hay tạo ra một phong ấn mới. Cách duy nhất để kết thúc là\dots\ \textit{`hoàn tất nghi lễ bị bỏ dở'}.''

``Hoàn tất nghi lễ?'' Azuki tròn mắt hỏi. ``Chẳng lẽ chúng ta phải tìm cách\dots\ đưa chúng xuống hố lại sao? Không thể nào!''

``Không, Azuki-san,'' Kaigou lên tiếng, giọng nói đều đều nhưng chắc nịch. Cô nhìn quanh một lượt. ``Không phải là lặp lại cái chết, mà là trao lại cho chúng những gì đã bị cướp đi. Một lời thú nhận, sự ghi nhận, và cho phép chúng được ra đi.''

``Chính xác,'' Shion mỉm cười đắc ý. ``Và để làm được điều đó, tài liệu này có ghi rõ, chúng ta cần năm vai trò thiết yếu cho một nghi lễ kết thúc. Một người mang \textit{`ký ức gốc'} để thấu hiểu nỗi đau. Một \textit{`thực thể trấn giữ'} để duy trì ranh giới giữa hai bờ sinh tử. Một \textit{`kẻ ngoại lai'} để phá vỡ cái vòng lặp do người trong làng tạo ra. Một người \textit{`kết nối'} có trái tim thuần khiết để giao tiếp với linh hồn. Và cuối cùng, một kẻ \textit{`đối diện trực diện'} để gánh vác áp lực và bảo vệ vòng tròn.''

Nghe đến đây, cả nhóm đưa mắt nhìn nhau. Một sự im lặng đồng thuận bao trùm căn phòng CLB chật hẹp. Không cần ai phải chỉ mặt đặt tên, cả năm chúng tôi đều hiểu rõ Shion đang ám chỉ những ai.

Người mang ký ức gốc, Misora Shino.

Thực thể trấn giữ, Shinomiya Sakura.

Kẻ ngoại lai phá vòng lặp, tôi, Hikawa Kuon.

Người kết nối thuần khiết, Hikawa Suzuki.

Và kẻ đối diện trực diện, vững chãi nhất, Hikawa Kaigou.

Năm mảnh ghép hoàn hảo đã được tập hợp, không phải do ngẫu nhiên, mà dường như là do sự sắp đặt của chính khu rừng Aogami. Đó cũng chính là năm người đã trực tiếp hồi sinh cả thị trấn này, không thẻ lẫn vào đâu được.

\pov{Hikawa Kaigou}

Sau khi năm mảnh ghép được điểm mặt gọi tên, bầu không khí trong phòng CLB bỗng chốc trở nên nghiêm túc hơn bao giờ hết. Chúng tôi bắt đầu bàn về kế hoạch tiến vào rừng Aogami một lần nữa.

``Nếu đã xác định được cách giải quyết, chúng ta không thể chần chừ,'' tôi khoanh tay trước ngực, nhìn thẳng vào cuốn tài liệu ố vàng. ``Càng để lâu, những linh hồn đó sẽ càng bức bối. Việc những đứa trẻ bắt đầu hiện hình rõ rệt ở trường học chứng tỏ ranh giới giữa hai thế giới đang mỏng đi rất nhiều.''

``Tôi đồng ý,'' Nanase lên tiếng, ánh mắt cô hiện lên sự quyết tâm. ``Dù không nằm trong năm vai trò chính, nhưng tôi và Mari-san sẽ đi cùng để hỗ trợ vòng ngoài. Ít nhất, với những gì đã trải qua, chúng tôi có thể giúp cảnh báo nếu có dị thường.''

``Tất nhiên là phải có tớ rồi!'' Ayame vỗ ngực cái \textit{bộp}. ``Hội học sinh không thể nhắm mắt làm ngơ khi học sinh của trường đang lao vào nguy hiểm được. Chút chuyện giữ trật tự và canh gác vòng ngoài cứ giao cho Onidoro Ayame này!''

``Cả mình nữa!'' Azuki giơ cao tay, dù vẻ mặt vẫn còn chút e dè. ``Mình cũng muốn đi. Từ trước đến giờ mình luôn là người dẫn đường, lần này mình muốn làm điều gì đó có ích hơn.''

Yuko-sensei thở dài, gõ nhẹ cuốn sổ điểm danh lên mặt bàn: ``Được rồi, nhưng với tư cách là giáo viên, cô bắt buộc phải đi theo để giám sát các em. Rừng Aogami vào ban đêm không phải là nơi để đi dã ngoại đâu.''

Tôi nhìn lướt qua nhóm hỗ trợ. Có họ đi cùng cũng tốt, ít nhất nếu có biến cố gì xảy ra ngoài dự kiến, Kuon và Suzuki sẽ không bị cô lập. Điều này sẽ đảm bảo những người vô tội sẽ không bị dính vào, vì điều đó sẽ không hay chút nào cả.

Tuy nhiên, câu hỏi quan trọng nhất vẫn còn bỏ ngỏ. Tôi lật lại tập tài liệu, mắt dừng ở đoạn chữ viết tay mờ nhạt bên lề một trang giấy.

``Shion-san,'' tôi ngẩng lên, hướng về phía cô gái tóc bạc-xám. ``Trong bản ghi chép có đề cập đến `chu trình thời gian' nhưng không nêu rõ ngày cụ thể. Cậu có thể đọc được khoảng thời gian nào thì năng lượng rừng đạt trạng thái cân bằng không?''

Shion khẽ gật, đôi mắt xanh đậm như hồ ngọc lặng gió. Cô đặt ngón tay lên một dòng chữ cổ, môi mấp máy đọc thầm rồi ngẩng lên. ``Tài liệu chỉ ghi `khi mặt trời thứ mười ba lần lặn sát ngưỡng cực'. Ta cũng không hiểu cho lắm đâu.''

Kuon đã mở điện thoại, gõ phím tính nhanh như chớp. ``Ngưỡng cực ở đây là điểm cực của đường hoàng đạo,'' cậu ta nhếch môi. ``Tính từ xuân phân, thứ mười ba lần mặt trời lặn sát ngưỡng cực rơi vào khoảng\dots\ từ 28 đến 30 tháng Năm.''

``Vậy là cuối tháng Năm này,'' tôi thở ra, cảm thấy không khí trong phòng như nặng thêm. ``Chúng ta có ba ngày đó thôi.''

``Vậy thì, Shino-san,'' tôi ngước về phía cô nàng tóc nâu đang đứng bên cạnh Kuon, ``theo cậu thì chúng ta sẽ đi vào thời điểm nào là thích hợp nhất?''

Shino ngước mắt lên, đôi đồng tử xanh đậm sâu thẳm ánh lên một thứ ánh sáng tĩnh lặng nhưng kiên định. Cô nhìn ra ngoài cửa sổ, nơi những tia nắng cuối ngày đang lụi tàn sau rặng thông già.

``Buổi tối, ba ngày nữa kể từ hôm nay,'' cô chậm rãi nói, giọng cô trong trẻo nhưng mang sức nặng của một lời phán truyền. ``Theo chu kỳ thời gian được ghi lại trong tài liệu, đó là thời điểm năng lượng của khu rừng chạm đến trạng thái cân bằng nhất. Khi đó, sự giao thoa sẽ giúp xoa dịu sự kích động của các linh hồn, tạo điều kiện tốt nhất để Suzuki-san kết nối với chúng.''

Cô quay lại nhìn tất cả chúng tôi, chốt hạ bằng một câu ngắn gọn:

``Chúng ta sẽ chấm dứt lời nguyền này. Một lần và mãi mãi.''

Tôi nhìn xuống bàn tay mình, nơi những đường gân xanh đang nổi rõ dưới làn da tái nhợt. Ba ngày, chỉ còn ba ngày nữa thôi. Mỗi nhịp đập của tim tôi như một tiếng trống thúc dồn, nhắc nhở rằng thời gian không chờ đợi ai. Những linh hồn kia đã chịu đựng hơn mười ba thế kỷ trong cô đơn, tủi nhục, và giờ đây, chúng đang trông chờ vào chúng tôi để kết thúc nỗi đau đó. Tôi không thể lùi bước. Dù có phải đối mặt với bóng tối tận cùng của lịch sử, dù có phải gánh lấy hậu quả nặng nề nhất, tôi cũng sẽ bước tới. Chúng tôi - bằng định mệnh của lịch sử - đã được chọn để đặt dấu chấm hết cho lời nguyền này, để những đứa trẻ kia được yên nghỉ, để Aogami không còn phải sống dưới cái bóng của quá khứ.

\begin{center}
  \uline{Ba ngày sau.}
\end{center}

Đúng bảy giờ tối, chúng tôi tập hợp trước bìa khu rừng Aogami, ở đúng vị trí mà chúng tôi đã đứng trong lần vào rừng trước.

Trời đêm không trăng, những tán thông đen kịt vươn lên giữa nền trời như những bàn tay khô khốc đang cố cào xé lớp màn nhung. Gió rít qua các kẽ lá tạo thành âm thanh rin rít, nghe vừa giống tiếng thở dài, lại vừa giống tiếng khóc nỉ non.

Trước khi chúng tôi bước qua ranh giới hàng rào gỗ mục, Kuon đột nhiên dừng lại. Cậu ấy tháo chiếc đồng hồ đeo tay với mặt kính kim loại lạnh lẽo ra, nắm lấy tay Suzuki và nhẹ nhàng đeo vào cho con bé.

``\vie{Suzuki, em giữ lấy cái này đi,}'' cậu nói, giọng trầm ấm nhưng ánh mắt vô cùng nghiêm túc.

Suzuki tròn mắt ngạc nhiên, ngón tay khẽ chạm vào mặt đồng hồ đang nhấp nháy một nhịp đèn đỏ li ti. ``\vie{Nhưng\dots\ Onii-sama, đây là đồ vật rất quan trọng của anh mà? Tại sao lại\dots}''

``\vie{Cứ nghe lời Kuon-kun đi, Suzu,}'' tôi vỗ nhẹ lên vai con bé. ``\vie{Em là người sẽ phải kết nối trực tiếp với bọn trẻ. Em cần một mỏ neo vững chắc hơn bất cứ ai trong chúng ta lúc này. Với cả, có lẽ em sẽ cần \emph{ông ta} hơn cả anh chị đấy.}''

Suzuki mím môi, rồi gật đầu cẩn thận vuốt lại quai đeo. ``\vie{Vâng, Onee-sama, Onii-sama. Em sẽ giữ nó cẩn thận ạ.}''

Tôi gật đầu, quay mặt nhìn về phía con đường mòn đang bị bóng tối nuốt chửng. Đây đã là lần thứ ba chúng tôi bước vào khu rừng này. Lần đầu tiên vào dạo trước, đó là một chuyến đi mù quáng, đầy rẫy sự tò mò và sợ hãi trước những hiện tượng siêu nhiên vô định. Lần thứ hai, sự thật phơi bày như một gáo nước lạnh, mang theo sự bàng hoàng và phẫn nộ tột cùng trước tội ác của lịch sử.

Nhưng lần này\dots\ không còn chỗ cho sự tò mò hay sợ hãi nữa. Chúng tôi mang theo một trọng trách nặng nề. Lần này, ngoài những chiếc đèn pin công suất lớn trên tay Azuki và Yuko-sensei, Shino còn cẩn thận mang theo một túi vải nhỏ chứa hạt lúa rang và một vài nhánh cỏ khô; chúng không phải để yểm bùa, mà để làm vật dẫn cho sự thanh tẩy, một cách để nói lời xin lỗi với mảnh đất này.

Đội hình được phân chia rõ ràng. Đi đầu là Shion và Azuki, họ sẽ dẫn đường quay lại đúng cái hố đất sâu thẳm mà chúng tôi đã tìm thấy trước đó. Theo ngay sau là nhóm hỗ trợ gồm Ayame, Nanase, Mari và Yuko-sensei. Shino và Sakura đi ở vị trí thứ ba để giữ ổn định tâm lý cho toàn đoàn. Và bộ ba Hikawa chúng tôi bọc lót ở phía sau cùng.

``Ái chà, cái rễ cây này mọc ngu ngốc thật đấy!'' tiếng Shion càu nhàu vang lên từ phía trước khi cô nàng suýt vấp ngã, làm ánh đèn pin loang loáng quét loạn xạ lên các thân cây.

``Suỵt! Shion-chan, cẩn thận dưới chân đi,'' Azuki khẽ nhắc nhở, tay cô giữ chặt lấy áo choàng của cô nàng Phù thủy. ``Đường đi trong này không giống như hành lang trường học đâu, nếu cậu ngã là không có ai đỡ được ngay đâu đấy.''

``Ta là đại pháp sư, làm sao có thể ngã vì dăm ba cái rễ cây phàm trần này được!'' Shion cãi lại, nhưng rõ ràng bước chân đã chậm và dè dặt hơn hẳn.

Phía sau họ vài bước, Sakura khẽ nghiêng đầu hỏi cô gái tóc nâu đang ôm khư khư túi vải nhỏ. ``Cậu ổn chứ, Shino? Năng lượng ở đây đang dao động rất mạnh, nếu thấy khó thở thì phải báo tôi ngay nhé.''

``Mình không sao, Sakura-san,'' Shino đáp, giọng nói điềm tĩnh đến lạ thường, như thể cô ấy đã hòa làm một với nhịp đập của khu rừng. ``Trái lại, mình cảm thấy chúng đang rất im lặng. Giống như\dots\ một lớp học đang nín thở chờ giáo viên bước vào vậy.''

``So sánh kiểu gì kỳ vậy trời\dots?'' cô nàng tóc hồng hơi nhăn mặt.

Tôi siết chặt nắm đấm, để móng tay bấm nhẹ vào lòng bàn tay nhằm giữ cho tâm trí hoàn toàn tỉnh táo. Dù Shino nói chúng đang im lặng, nhưng cái lạnh buốt xương của Aogami vẫn không ngừng luồn lách qua lớp áo khoác.

``Đi sát vào nhau,'' Kuon nhắc nhở tôi và Suzuki bằng giọng đủ nghe. ``Dù có nghe thấy hay nhìn thấy gì, tuyệt đối không được tự ý tách đoàn.''

``Vâng!'' Cả tôi và Suzu đồng thanh.

Chúng tôi tiếp tục tiếp bước, vừa mò mẫm dò đường, vừa thi thoảng nói chuyện để làm giảm bầu không khí căng thẳng và hồi hộp. Càng đi sâu vào rừng, sự khác biệt so với hai lần trước càng hiện ra rõ rệt.

Lần đầu tiên, con đường phủ đầy lá mục và tĩnh lặng đến gai người. Lần thứ hai, cái lạnh của thời gian và những bóng người thời cổ đại ám lấy từng ngọn cây. Nhưng đêm nay, khu rừng Aogami không im lìm, cũng chẳng tua lại quá khứ. Nó đang \textit{thức giấc}.

Thay vì bóng tối đặc quánh, một lớp sương mù mỏng, màu xám lấp lánh như lân tinh bắt đầu bò ra từ dưới những lớp rêu phong. Nó trườn qua mắt cá chân chúng tôi, mang theo cái lạnh buốt thấu xương. Nhờ ánh đèn pin của Azuki quét qua, tôi thề rằng mình đã nhìn thấy vô số vết chân nhỏ xíu - cỡ của trẻ lên năm, lên bảy - in hằn trên nền đất ẩm, rồi tan biến ngay tức khắc.

Không gian đặc nghẹt những âm thanh rì rầm. Không phải là gió rít qua tán lá nữa, mà là những tiếng cười rúc rích khe khẽ, xen lẫn tiếng sụt sùi đứt quãng. Thỉnh thoảng, từ phía sau những gốc thông đen khổng lồ, một vài bóng trắng nhỏ xíu ló ra, đôi mắt rỗng không nhìn chằm chằm vào đoàn người rồi rụt lại khi ánh sáng đèn pin vô tình lướt tới.

``Mọi người\dots'' giọng Mari-san run lên bần bật từ phía giữa hàng. ``Có\dots\ có ai vừa kéo vạt áo mình không?''

``Đừng quay lại nhìn,'' Yuko-sensei cất giọng rành rọt, dù tay cô đang nắm chặt lấy bả vai của cô nàng tóc đỏ-nâu. ``Cứ đi thẳng đi. Chúng chỉ đang thử sự chú ý của chúng ta thôi.''

Nanase rút trong túi áo ra một chuỗi hạt gỗ, miệng lẩm nhẩm những câu cầu nguyện đứt quãng. Ngay cả Ayame, người luôn tự hào về sự bạo dạn của mình, lúc này cũng im bặt, hai tay khoanh chặt trước ngực như đang cố tự sưởi ấm.

Sự hiện diện của các linh hồn đêm nay quá đông đúc. Chúng không tấn công vật lý, nhưng sức ép tâm lý mà chúng tạo ra nặng nề hệt như một lớp đất đang từ từ đè lên ngực, ép hết không khí ra ngoài. Tôi có thể nghe rõ tiếng thở dốc của Suzu bên cạnh, và bàn tay con bé đang bám chặt lấy ống tay áo của Kuon đến mức các khớp ngón tay trắng bệch.

Dù vậy, chiếc đồng hồ trên cổ tay Suzu vẫn đang phát ra từng nhịp chớp đỏ đều đặn, như một nhịp tim cơ học chống lại sự nhiễu loạn của thế giới tâm linh. Nó thực sự đang giúp con bé giữ vững mỏ neo của mình.

``Sắp đến nơi rồi,'' tiếng Shion vang lên từ phía trước, giọng cô nàng có phần nghèn nghẹt chứ không còn sự ngạo nghễ như lúc nãy. ``Cái hố\dots\ nó nằm ngay sau gò đất kia.''

\pov{Hikawa Suzuki}

Đúng lúc Shion-san vừa dứt lời, bàn chân tôi bất ngờ vướng phải một đoạn rễ thông trồi lên khỏi mặt đất. Tôi loạng choạng và ngã khuỵu xuống lớp lá mục ẩm ướt.

``Em ổn chứ, Suzu?'' Tiếng Onee-sama vang lên từ phía trước.

Tôi vội xua tay, cố gắng gỡ chiếc dây giày đang bị vướng chặt vào mấu gỗ sần sùi. ``Em không sao ạ! Mọi người cứ đi trước đi, em gỡ xong sẽ chạy theo bắt kịp ngay thôi.''

``Nhanh lên đấy,'' Onii-sama nhắc nhở rồi tiếp tục bước theo ánh đèn pin đang dần khuất sau gò đất.

Chỉ vài giây ngắn ngủi sau khi ánh sáng cuối cùng của đoàn người biến mất khỏi tầm mắt, bóng tối đã lập tức ùa tới, đặc quánh và ngột ngạt. Bất chợt, mặt đồng hồ trên tay tôi chớp đỏ liên tục. Một giọng nói đặc sệt âm hưởng Pháp vang lên từ chiếc loa nhỏ, cất lên bằng tiếng Anh:

``\eng{\dots Cô biết không, Suzuki-user? Trong mọi bộ phim kinh dị rẻ tiền mà ta từng xem, kẻ nào ngu ngốc tự ý tách khỏi đoàn, kẻ đó chắc chắn sẽ là người đầu tiên phải bỏ mạng đấy.}''

Tôi khẽ rùng mình, cố lờ đi lời điềm báo rợn tóc gáy của ông Game để tập trung tháo sợi dây giày ngoan cố. ``\eng{Đ-Đừng làm cháu sợ chứ, G-Game-san\dots}''

``\eng{Eh, chỉ là suy nghĩ vẩn vơ của tôi thôi. Trường hợp này khác kha khá so với cái mô-típ vớ vẩn kia đấy,}'' ông ấy đáp lại bằng giọng có gì đó vẫn bỡn cợt. Hoàn toàn chẳng hề giúp tôi bình tĩnh hơn được cả. Tôi cố gắng hít một hơi thật sâu, rồi nhanh chóng luồn tay xuống dưới gấu quần để tháo gỡ nút thắt phức tạp.

Nhưng ngay khi tôi vừa rút được chân ra, những tiếng bước chân lạo xạo đạp lên lá khô vang lên từ phía sau lưng. Không phải là những bước chân nhỏ bé của các linh hồn, mà là bước chân nặng nề, thô kệch của con người.

Tôi quay ngoắt lại. Dưới ánh trăng lờ mờ xuyên qua tán lá, bốn cái bóng cao lớn hiện ra.

Đi đầu là gã trùm côn đồ với mái tóc nhuộm vàng chóe, kẻ đã từng gây rắc rối với Onii-sama trong ngày thứ hai anh học. Hai tên đàn em thân tín của gã đứng hai bên, trên tay lăm lăm mấy khúc gậy gỗ. Đi tụt lại phía sau một chút là một cậu học sinh nam trông khá quen mắt, dường như là nạn nhân thường xuyên bị bọn chúng sai vặt mua đồ ở trường. Trông cậu ta rụt rè, tay ôm chặt lấy balo, mặt cúi gằm xuống đất như bị ép buộc phải tới chốn ma quỷ này.

Bọn chúng hoàn toàn không biết chúng tôi đang ở đây để thực hiện một nghi thức nguy hiểm.

``Ái chà chà, nhìn xem chúng ta tìm được ai này?'' Gã tóc vàng nhếch mép cười cợt, nhổ toẹt bã kẹo cao su xuống đất. ``Kia chẳng phải là cô em gái bé nhỏ ngoan ngoãn của con ranh Kaigou sao?''

``Nửa đêm nửa hôm vác xác vào rừng một mình, chắc là đang cô đơn lắm đây,'' một tên đàn em hùa theo, cười hô hố bằng chất giọng khả ố.

Tôi lùi lại một bước, trái tim bắt đầu đập liên hồi trong lồng ngực. ``Các người\dots\ muốn gì?''

``Muốn gì à? Chỉ muốn mượn em gái đây một chút để `nhắn gửi' vài lời yêu thương tới thằng anh trai và con chị gái xấc xược của cưng thôi!''

Nói rồi, chưa để tôi kịp phản ứng, gã tóc vàng lao tới như một con thú đói. Bàn tay thô ráp của gã siết chặt lấy cổ tay tôi, kéo giật mạnh khiến tôi ngã nhào về phía trước. Tôi toan há miệng hét lên gọi Onii-sama, nhưng ngay lập tức, tên đàn em còn lại đã chồm tới, dùng một bàn tay bẩn thỉu bịt chặt lấy miệng tôi.

``Ưm\dots! Ưm\dots!'' tôi giãy giụa tuyệt vọng, nhưng sức lực của một đứa con gái hoàn toàn vô vọng trước ba gã côn đồ to xác. Chúng thô bạo kéo lê tôi xềnh xệch trên nền lá mục, rẽ sang một nhánh đường hoàn toàn khác với hướng mà nhóm Onii-sama vừa đi.

Trong lúc bị kéo ngang qua cậu nam sinh đi cùng chúng, tôi thấy cậu không hề hùa theo cười cợt. Cậu ấy chỉ đứng run rẩy, ánh mắt lảng tránh cái nhìn cầu cứu của tôi, miệng lẩm bẩm đủ để tôi nghe thấy:

``Xin lỗi cậu nhé\dots\ Suzuki-san.''

Nước mắt tôi trào ra, nóng hổi và bất lực. Khu rừng xung quanh vẫn im lìm, nhưng tôi biết, trong bóng tối kia, những linh hồn không hề bỏ qua cảnh tượng này.

\pov{Hikawa Kaigou}

Sau một hồi lội qua lớp sương mù đặc quánh và những bụi gai cào rách cả gấu quần, chúng tôi cuối cùng cũng đến được trung tâm của khu vực cấm. Trước mắt chúng tôi, dưới ánh sáng leo lét của đèn pin, là miệng hố đất rộng hoác, nơi đã từng nuốt chửng không biết bao nhiêu sinh mạng vô tội.

``Trời đất\dots'' Ayame lùi lại một bước, tay bưng lấy miệng. Sự bạo dạn và vẻ kiêu hãnh thường ngày của Hội trưởng Hội học sinh bay biến sạch. ``Đây\dots\ đây là nơi đó sao? Rộng đến mức này ư?''

``Chỉ nhìn thôi cũng cảm thấy rùng mình,'' Nanase ôm lấy hai vai, giọng cô run lên. ``Họ thực sự đã bắt những đứa trẻ bước xuống cái hố sâu thẳm này sao? Thật khủng khiếp\dots''

Azuki đứng sát bên tôi, môi cô mím chặt đến mức tái nhợt. Cô gái dẫn đường thường ngày luôn tự tin giờ đang run nhè nhẹ, tay siết chặt đèn pin đến mức các khớp trắng bệch. ``Lần trước\dots\ lần trước mình đã nghĩ mình chuẩn bị kỹ rồi,'' cô thì thầm, giọng khàn đặc, ``nhưng mà nhìn thấy lại cái hố này\dots\ nó vẫn như cái miệng quỷ thật sự.''

Mari ôm chặt cuốn sổ vẽ vào ngực, ngón tay cô run rẩy đến mức từng trang giấy lật đánh rột rột. Cô không nói gì, chỉ gật nhẹ, mắt vẫn dán vào miệng hố như thể chỉ cần chớp mắt một cái, những bàn tay trẻ con sẽ túm lấy cô kéo xuống. Một giọt mồ hôi lạnh chảy dài trên thái dương cô, ướt cả vành tai.

Yuko-sensei cũng không giấu được vẻ căng cứng dù đây là lần thứ hai mà cô chiêm ngưỡng. Cô đứng thẳng, nhưng bả vai hơi co lên, hai tay khoanh trước bụng theo phản xạ tự vệ. Đôi mắt nghiêm nghị sau cặp kính lấp lánh ánh đèn pin, nhưng khóe miệng cô giật nhẹ. ``Đừng để cảm xúc lấn át,'' cô ra lệnh, giọng trầm hơn bình thường một nốt, ``nhưng\dots\ lần này nếu có gì bất thường, tôi sẽ cho cả đoàn rút lui ngay lập tức.''

Trong khi nhóm hỗ trợ vẫn còn đang bàng hoàng, lớp sương mù dưới đáy hố đột nhiên cuộn lên. Một ánh sáng màu xanh lục lờ mờ tỏa ra, và từ trong đó, một hình bóng nhỏ nhắn từ từ bước tới. Đó là cô bé ấy, với mái tóc tết đuôi sam và con búp bê bằng vỏ sò và rơm rạ ôm chặt trong tay. Đôi mắt xanh lục của em nhìn lướt qua chúng tôi. Không còn u uất hay mệt mỏi như lần chạm mặt trước ở chuyến đi lần hai, lần này, ánh mắt em điềm tĩnh đến lạ thường.

``Vậy là mọi người đã tới thật,'' cô bé cất tiếng, giọng mỏng manh nhưng rõ ràng, vang vọng giữa bốn bề tĩnh lặng.

Shino bước lên một bước, đối diện với cô bé. Cô khẽ gật đầu, ánh mắt dịu dàng nhưng đầy kiên định: ``Chị đã hứa mà. Chị mang ký ức của các em, và chị sẽ không để nó bị chôn vùi thêm một lần nào nữa.''

Tôi siết chặt tay thành nắm đấm, nói thêm: ``Bọn chị ở đây để sửa chữa những gì lịch sử đã làm sai. Không phải bằng cách lấp liếm cho tội ác, mà bằng cách đối diện với nó.''

``Bọn anh sẽ đưa các em đến nơi mà các em thuộc về,'' Kuon tiến lên, đứng cạnh Shino. ``Một nơi an nghỉ thực sự.''

Tôi nhìn vào đôi mắt xanh lục kia và bất chợt hiểu: nỗi đau lớn nhất của em không phải là cái chết, mà là bị từ chối quyền được khóc, được sợ, được yêu thương, bị biến thành một biểu tượng không ai công nhận thay vì một đứa trẻ. Trong khoảnh khắc ấy, tôi cảm thấy dòng máu mình như chảy ngược; tôi cũng từng là đứa trẻ bị cả lớp nhìn qua, chỉ vì mang họ Hikawa, chỉ vì không thể cười thật to. Sự công nhận mà chúng tôi đem đến hôm nay không chỉ là lời xin lỗi, mà là một tấm gương soi lại chính mình: nếu không dám nhìn thẳng vào nỗi đau của người khác, ta sẽ mãi mãi để lại phía sau những linh hồn lạc lõng, kể cả linh hồn của chính ta.

Cô bé mắt xanh khẽ nghiêng đầu, đôi môi mím lại như đang cố kiềm nén một cảm xúc nào đó đã bị đè nén suốt mười ba thế kỷ. Em chậm rãi gật đầu, lùi lại một bước để nhường khoảng trống cho chúng tôi bắt đầu.

``Được rồi, vào vị trí thôi,'' Shion hắng giọng, mở tập tài liệu cổ ra. ``Thời khắc giao thoa đã đến. Chúng ta phải bắt đầu ngay thôi-- Ơ kìa?''

Tên Phù thủy tự xưng dừng lại, mắt mở to nhìn quanh: ``Suzuki đâu rồi?''

``\dots Hả?''

Câu hỏi của Shion như một gáo nước lạnh tạt thẳng vào mặt tôi. Tôi ngoắt đầu lại, nhìn lướt qua Ayame, Nanase, Mari, Azuki, và Yuko-sensei.  Không có bóng dáng mái tóc đen cột hai bím rụt rè quen thuộc đó đâu cả.

``Suzu? \textbf{Suzu!}'' tôi hoảng hốt gọi lớn, toan quay đầu chạy ngược lại con đường mòn. ``Thôi chết rồi, không lẽ nào con bé lại--''

Nhưng Azuki đã đưa cánh tay ra chặn tôi lại. ``Bình tĩnh đi, Kaigou-chan.''

``\textbf{Bình tĩnh thế nào được!}'' tôi gắt lên, tim đập thình thịch trong lồng ngực. ``Con bé bị tụt lại rồi! Giữa cái khu rừng đầy rẫy oán khí này!''

``Azuki-san nói đúng đấy, Kaigou-chan,'' Kuon bước tới, đặt tay lên vai tôi trấn an. ``Cô không nhớ sao? Đây đâu phải lần đầu hai chị em cô bị tách nhau ra. Lần ở The Void, hay lần mà em ấy chạy đi tìm Shino-san, Suzuki đều đã tự mình xoay sở được cơ mà.''

``Cậu nên tin tưởng vào em gái của mình đi,'' Azuki nhìn thẳng vào mắt tôi, giọng nói chắc nịch. ``Cậu ấy trông thế thôi mà rất mạnh mẽ và tự tin lắm đấy chứ. Cậu ấy cũng không còn bé bỏng để cậu có thể bao bọc được nữa đâu.''

Yuko-sensei gật đầu đồng tình: ``Tôi sẽ nhờ Akagane-san và Furukawa-san đứng cảnh giới phía sau, nếu có động tĩnh gì từ Suzuki-san sẽ lập tức báo lại. Còn bây giờ, nghi lễ là ưu tiên hàng đầu.''

Tôi định lên tiếng phản bác, thì Shion đã lật tức thúc giục. ``Không thể chờ được đâu!'' Cô sốt sắng lật trang giấy. ``Thời điểm linh thiêng nhất chỉ kéo dài một khoảng ngắn thôi, nếu lỡ mất, năng lượng sẽ bị đảo lộn và tỷ lệ thành công sẽ về không!''

Tôi cắn chặt môi, cố nén lại cơn hoảng loạn đang cào xé ruột gan. Lời của Kuon và Azuki dần kéo tôi về với thực tại. Phải, Suzu mạnh mẽ hơn tôi tưởng, kể từ khi bị lạc ở The Void, băng qua trong thế giới Cầu Vồng. Và con bé đang giữ chiếc đồng hồ của Kuon nữa. Nó biết là nó sẽ phải làm gì.

Tôi hít một hơi thật sâu, ép mình phải bình tĩnh. Tôi hỏi cô nàng Phù thủy tự xưng kia: ``Trình tự nghi thức thế nào, Shion-san? Vai trò của `Người kết nối' nằm ở đâu?''

Shion liếc nhanh qua tập tài liệu: ``Nó nằm ở bước cuối cùng! `Người kết nối' sẽ tạo cầu nối cuối cùng để giải thoát các linh hồn sau khi ranh giới đã được mở và ký ức được gọi về.''

Nghĩa là chúng tôi có thể bắt đầu trước mà không cần Suzu ngay lúc này. Tôi gật đầu, lùi lại vị trí của mình - đóng vai trò là ``Kẻ đối diện trực diện''.

``Được,'' tôi hít sâu một lần nữa, hướng mắt về phía miệng hố. ``Chúng ta sẽ bắt đầu nghi thức.''

Nhưng trong vang vảng suy nghĩ tôi, tôi vẫn không tài nào không cảm thấy lo lắng cho cô em gái của mình. ``\textit{Suzu, mong là em có thể tới kịp\dots}''

\pov{Hikawa Suzuki}

Bọn chúng lôi tôi vào một khoảng trống giữa bụi cây gai, nơi ánh đèn pin từ đoàn đã khuất hoàn toàn. Một tên rút sợi dây thừng xanh dài từ trong ba lô, quấn vòng quanh thân tôi từng vòng siết chặt. Dây cắt sâu vào làn da qua lớp áo khoác mỏng, khiến tôi không thể cử động khuỷu tay hay đầu gối. Gã tóc vàng giật mạnh nút thắt, kéo tôi ngã nghiêng dựa vào gốc thông già sần sùi. Vỏ cây lởm chởm cào vào lưng, hơi ẩm của đất mục bốc lên nồng nặc. Mùi lá mục, thuốc lá và mồ hôi nam tính hòa quyện khiến tôi nghẹt thở. Bốn bề là bóng cây cao vút, tán lá rậm rạp che kín mặt trăng, chỉ để lại những khe hở nhỏ xíu lấp lánh như mắt quỷ. Tiếng gió rít qua kẽ lá nghe như tiếng thì thầm chế giễu. Tôi cố giật mình, nhưng dây thừng chỉ siết chặt hơn, đẩy đầu tôi cúi gằm xuống đám rêu ướt. Chiếc đồng hồ của Onii-sama vẫn đập nhịp đỏ đều đặn trên cổ tay bị trói, như trái tim cô độc giữa rừng đêm.

Nhưng cảm giác này\dots\ tôi đã từng trải qua rồi. Không phải ở khu rừng Aogami đầy ma quái này, mà là ở thế giới cũ với Onee-sama, nơi tôi từng là một đứa trẻ lạc lõng, mang trong mình dòng máu Hikawa bị người đời khinh miệt. Những kẻ như gã tóc vàng này, tôi đã gặp đủ. Họ không cần lý do để làm tổn thương người khác; họ chỉ cần biết rằng tôi yếu thế, rằng tôi không có chỗ dựa. Tôi từng bị giữ lại sau giờ học, bị đẩy vào nhà vệ sinh nữ và bị dọa ``dạy dỗ'' chỉ vì tôi không cười khi chúng chế nhạo thầy giáo. Tôi từng bị lột sách vở, bị quay clip rồi tung lên mạng với dòng chữ: ``Con nhãi Hikawa này mà cũng học giỏi à?'' Và khi tôi khóc, chúng cười. Khi tôi chạy, chúng đuổi. Giống như đêm nay. Giống hệt như đêm nay.

Chỉ khác một điều: lần này, tôi không còn là đứa trẻ chỉ biết ôm mặt khóc trong góc lớp. Tôi sẽ không để chúng nó được thoả mãn đâu.

``Này, đại ca, con nhỏ này còn run cầm cập nữa kìa!'' tên đàn em bên trái cười khẩy, dùng đầu ngón chân đạp nhẹ vào vai tôi. ``Có khi nó đang tưởng tượng ra mấy con ma trong rừng nên sợ đái ra quần rồi!''

``Đừng có đụng vào tôi!'' tôi ngẩng đầu lên, cố giữ giọng không run. ``Các người không biết mình đang đứng ở đâu\dots\ Nơi này không phải chỗ để đùa giỡn!''

Tên còn lại nhếch mép, phì cười: ``Đại ca ơi, nó còn dám dọa bọn mình kìa! Chắc nó tưởng mình là thánh nữ hay gì đó, định dùng lời thánh thót để cứu rỗi bọn mình đấy nhể!''

Gã tóc vàng bước lại gần, nắm tóc tôi giật mạnh khiến đầu tôi ngửa ra sau. Gã hướng gưong mặt hầm hố ra gần tới chỗ tôi, khịt mũi: ``Mày nghĩ bọn tao không dám làm gì mày sao? Một con nhãi như mày mà dám nói mấy lời đao to búa lớn đấy hử?''

Một tên đàn em khác móc trong túi ra chiếc điện thoại, bật camera nháy sáng. ``Đại ca, hay là ta quay lại cảnh con nhỏ này khóc lóc lột đồ đi. Chốc nữa gửi cho thằng anh nó, bắt nó phải lết đến đây quỳ lạy liếm giày cho đại ca! Không thì sáng mai cả cái trường này sẽ có phim hay để xem.''

``Ồ\~{} không tồi đâu. Tao thích ý tưởng đó đấy.''

Nghe đến đó, mặt tôi tái mét. Nỗi sợ hãi trào dâng, nhưng cùng với đó là một sự phẫn nộ không thể kìm nén. Bọn chúng không chỉ muốn đánh tôi, chúng muốn hủy hoại danh dự của gia đình Hikawa. Tôi gồng mình, dùng hết sức bình sinh cắn ngập răng vào bàn tay đang bịt miệng tôi.

``\textbf{Á!} Con chó này!'' tên đàn em rú lên thảm thiết, vội vã rụt tay lại.

Nhân cơ hội đó, tôi dùng cùi chỏ giáng mạnh vào bụng tên đang giữ tay mình và hét lớn: ``\textbf{Cứu em với! Onii-sama--}''

Chưa kịp dứt lời, gã tóc vàng đã túm chặt lấy cổ áo tôi, ném mạnh tôi về phía sau. Lưng tôi va đập vào gốc thông sần sùi đau điếng, tầm nhìn tối sầm lại. Trong ánh mắt mờ mờ do tác động mạnh ấy, tôi thấy gã tiến lại gần, cười khẩy lần nữa, coi tôi như mấy trò tiêu khiển: ``Mày gan nhỉ? Còn có sức chống trả bọn này đấy? Có giỏi thì hét to lên xem, tao cho mày đấm vỡ a lô bây giờ!''

Nhưng trước khi gã kịp lao tới giở trò tiếp, chiếc đồng hồ trên tay tôi đột ngột phát ra một tiếng \textit{tít} dài chói tai. Mặt kính kim loại lóe lên một luồng ánh sáng đỏ rực rỡ, chói lòa như con mắt ác quỷ trong đêm tối. Và rồi, một giọng nói lạnh lẽo, vô cơ nhưng mang đầy sự khinh bỉ vang lên bằng tiếng Nhật chuẩn xác: ``Chậc chậc. Nhắm vào mắt xích yếu nhất trong đội hình để trả thù cá nhân. Một nước đi quá sức đê hèn và hạ đẳng, hệt như cái não phẳng lặng của các ngươi vậy. À mà không, gọi các ngươi là rác rưởi có lẽ còn xúc phạm đến cả khái niệm của rác rưởi.''

Ba tên côn đồ giật thót mình, dáo dác nhìn quanh: ``\textbf{Ai? Thằng nào vừa nói đó? Ra mặt đi!}''

Cậu nam sinh đứng phía sau lùi lại vài bước, mặt cắt không còn một giọt máu, chỉ tay bằng cánh tay run rẩy: ``Đ-Đại ca\dots\ g-giọng nói đó\dots\ phát ra từ cái đồng hồ của cậu ấy ạ\dots''

``Đồng hồ cái chó gì!'' gã tóc vàng nổi điên, rút từ trong túi quần ra một con dao bấm kêu lách cách. Hắn nhìn tôi với đôi mắt vằn lên những tia máu đỏ lừ và chĩa nó vào tôi: ``Mày giấu cái loa gì trong đó hả con ranh? Dám giở trò với tao à? Hôm nay tao sẽ rạch nát cái mặt mày!''

``\textbf{Dừng lại!}'' tôi hét lên, cố gắng chống tay đứng dậy, nhưng mắt cá chân trái đã truyền đến cơn đau nhói do cú ngã ban nãy. Dù vậy, tôi nhất quyết không lùi bước. Tôi trừng mắt nhìn thẳng vào gã, ánh mắt kiên cường mà tôi học được từ chị tôi: ``\textbf{Các người không biết mình đang đứng ở đâu đâu! Chỗ này đang tràn ngập--}''

``\textbf{Tao đếch quan tâm!}'' hắn gầm lên, giơ chiếc giày bọc mũi thép lên cao, chuẩn bị giáng một cú đạp toàn lực vào bụng tôi.

Nhưng trước khi đế giày của hắn kịp chạm vào người tôi, một âm thanh trầm đục vang lên từ dưới lòng đất. Không phải là động đất, mà giống như một nhịp tim khổng lồ vừa đập mạnh một nhịp. Toàn bộ khu rừng như rùng mình. Những tán thông đen kịt rung bần bật, trút xuống một cơn mưa kim thông sắc nhọn.

``Ngu ngốc,'' giọng nói của Game lại vang lên, lần này trầm đục và mang âm hưởng của một thực thể siêu việt. ``Ta đã nói đây là một bối cảnh phim kinh dị. Và trong những bối cảnh như thế này, những kẻ cặn bã bắt nạt kẻ yếu\dots\ luôn nhận một kết cục thảm khốc nhất.''

Nhiệt độ xung quanh giảm sút đột ngột, lạnh buốt đến mức hơi thở của chúng tôi lập tức hóa thành khói trắng. Lớp sương mù lấp lánh dưới chân tôi không còn là những vệt mờ nhạt nữa. Chúng cuộn lên, xé toạc không gian, đặc quánh và tanh nồng mùi rêu mốc. Từ trong màn sương ấy, vô số đôi tay nhỏ xíu với làn da nhợt nhạt, chi chít vết xước vươn ra, bám chặt lấy mắt cá chân của ba gã côn đồ.

Và rồi, những bóng ma trẻ em hiện hình. Chúng đứng vòng quanh, đôi mắt rỗng không nhìn chằm chằm vào những kẻ vừa buông lời nhục mạ.

``Cái\dots\ cái quái gì thế này?!'' tên cầm điện thoại đánh rơi chiếc máy xuống đất khi thấy màn hình chớp nháy loạn xạ rồi hiện ra một khuôn mặt trẻ con đẫm máu đang từ từ bò ra khỏi màn hình.

``Thả tao ra! B-Bớ người ta, ma!! \textbf{Ma! Có ma!}'' gã tóc vàng điên cuồng vung dao chém loạn xạ vào không khí, nhưng những đôi tay ma quái đó không hề hấn gì.

Cái giá lạnh của Aogami bắt đầu đánh gục tâm trí của chúng. Ảo ảnh kinh hoàng ập tới. Qua tiếng gào thét của chúng, tôi có thể lờ mờ nhận ra chúng đang nhìn thấy gì: bóng tối của huyệt mộ, cảm giác nghẹt thở khi từng xẻng đất lạnh lẽo đổ ập lên người, và sự tuyệt vọng của những đứa trẻ bị chôn sống. Gã trùm ôm lấy cổ, há hốc miệng cố gắng hớp lấy không khí như thể đang bị vùi lấp dưới lớp đất đá vô hình. Hai tên đàn em thì cào cấu gầm gừ trên nền lá mục, cố gắng đào bới một lối thoát khỏi cái hố tưởng tượng. Một cảnh tượng mà những linh hồn trẻ thơ ấy đã từng trải qua, hết lớp này tới lớp khác, và giờ bọn chúng đang trải nghiệm lại đúng cái cảm giác đó.

Giữa sự hỗn loạn và gào thét kinh hoàng đó, cậu học sinh nam - người nãy giờ vẫn đứng rụt rè phía sau và dường như là kẻ duy nhất không bị đám linh hồn nhắm tới - đã lấy hết can đảm lao tới. Cậu dùng một mảnh đá nhọn nhặt được gần đó, run rẩy nhưng dứt khoát cứa đứt sợi dây thừng đang trói nghiến lấy tôi.

``Suzuki-san, cậu mau chạy đi!'' cậu ấy thì thầm, đỡ tôi đứng dậy. ``Tớ xin lỗi vì đã không giúp cậu sớm hơn\dots''

Tôi gật đầu hàm ơn. Sau cú ngã ban nãy và cú ném của tên trùm kia, toàn thân tôi ê ẩm, đặc biệt là ở cổ chân bị trói, giờ lại càng đau nhói hơn mỗi khi cử động. Tuy nhiên, tôi biết mình không có thời gian để nghỉ ngơi hay than vãn nữa. Mọi người vẫn đang chờ tôi.

Tôi vừa kịp đứng vững thì sương mù xung quanh bắt đầu chuyển động. Lớp sương không còn xoáy quanh bọn côn đồ nữa mà từ từ tản ra, tạo thành một con đường hẹp phát sáng mờ ảo dẫn sâu vào trong rừng, đúng hướng mà nhóm Onii-sama vừa đi. Những bóng ma trẻ con cũng ngừng việc hù dọa đám người kia, chúng quay lưng lại, chậm rãi lướt đi trên con đường sương mù đó, thỉnh thoảng lại ngoái đầu nhìn tôi như thúc giục.

``\eng{Đi theo chúng đi,}'' giọng nói của Game từ chiếc đồng hồ vang lên đều đều, bình thản. ``\eng{Chúng đang dẫn đường cho cô đến nơi tiến hành nghi lễ đấy, Suzuki-user. Thời khắc giao thoa sắp bắt đầu rồi.}''

Tôi không chần chừ thêm nữa, cắn răng phớt lờ cơn đau ở cổ chân, khập khiễng bước theo những linh hồn nhỏ bé.

Ở phía sau, khi áp lực tâm linh của khu rừng giảm bớt, ba tên côn đồ dần lấy lại được chút tỉnh táo. Chúng lồm cồm bò dậy, quần áo xộc xệch, dính đầy đất bùn. Gã tóc vàng thở hổn hển, đôi mắt vằn tia máu nhìn về phía tôi đang xa dần.

``Mẹ kiếp\dots\ con ranh đó\dots\ nó chơi bùa ngải!'' gã rít lên qua kẽ răng, vung tay tát mạnh vào mặt cậu nam sinh vừa cởi trói cho tôi. ``\textbf{Mày đứng đực ra đó làm đếch gì!? Đuổi theo nó! Đừng hòng tao tha cho nó đêm nay!}''

Bất chấp nỗi sợ hãi tột cùng vừa trải qua, sự tức giận và thói hung hăng mù quáng đã lấn át lý trí của bọn chúng. Chúng lảo đảo chạy đuổi theo con đường sương mù, và dĩ nhiên, cậu nam sinh rụt rè kia lại một lần nữa bị kéo đi một cách ép buộc.

\pov{Hikawa Kaigou}

Tại miệng hố trung tâm, sương mù bắt đầu cuộn xoáy dữ dội thành những luồng khí lạnh ngắt, bắt đầu cho nghi lễ siêu thoát.

``\textbf{Bắt đầu đi!}'' Shion hô lớn, mở bung tập tài liệu cổ ra trước mặt. ``Sakura-san, thiết lập ranh giới!''

Sakura lập tức rút từ trong túi áo ra ba tấm bùa giấy nhạt màu, ném chúng về ba hướng xung quanh miệng hố. Cô khép mắt, thì thầm từng chữ rõ ràng: ``Khai môn, tĩnh giới, cầu hộ thân.'' Sau đó, cô đập mạnh hai bàn tay vào nhau, tạo ra một tiếng vang khô khốc. Những tấm bùa bốc cháy, tạo ra một bức tường năng lượng mỏng manh, ánh lên màu vàng nhạt, ngăn không cho sương mù đen từ bên dưới hố tràn ra ngoài.

``Đã cố định ranh giới!'' cô nàng Vu nữ nói lớn, trán bắt đầu lấm tấm mồ hôi. ``Shino-san, tới lượt cậu!''

Shino bước lên trước, đối diện với cô bé mắt xanh đang lơ lửng lặng im giữa không trung. Cô lấy từ trong túi vải ra những hạt lúa rang và nhánh cỏ khô, thả từ từ xuống lòng hố. Mỗi hạt lúa chạm đất, một âm thanh vang lên như tiếng chuông gió.

``Chị mang ký ức của các em,'' Shino nhắm mắt lại, giọng nói trong trẻo vang vọng khắp khu rừng. ``Ký ức về những buổi chiều chơi đùa, về những bài hát đồng dao, và cả nỗi sợ hãi khi bị bóng tối nuốt chửng. Chị nhớ tất cả. Và chị gọi tên các em, để các em biết rằng mình chưa bao giờ bị lãng quên.''

Dưới lòng hố, những đốm sáng nhỏ li ti bắt đầu bay lên, bám lấy những hạt lúa rang. Cô bé mắt xanh ngước nhìn Shino, đôi môi khẽ mở: ``Tại sao\dots\ bọn em lại bị bỏ lại?''

``Em biết không? Những người lớn khi ấy đã quá yếu đuối và ích kỷ,'' Kuon đáp, giọng điềm tĩnh nhưng kiên quyết. ``Họ đã chọn cách trốn tránh thay vì bảo vệ các em. Có thể anh chị chắc chắn không thể thay đổi quá khứ, nhưng mọi người đang ở đây để thừa nhận tội ác đó. Các em không có lỗi. Các em không đáng bị bỏ lại.''

``Các em xứng đáng được nghe lời xin lỗi,'' Azuki cất tiếng, mắt long lanh. ``Và xứng đáng được đưa về nhà.''

Cô bé mắt xanh ngẩng lên, đôi mắt trong veo như chứa cả bầu trời đêm Aogami. Giọng em nhỏ nhẹ, run run như sợ gió tanh sẽ cuốn mất lời: ``Các anh chị\dots\ có thể đưa bọn em quay trở về không? Về nơi\dots\ có mẹ đang đợi, có bạn bè đang chơi trốn tìm dưới ánh nắng\dots\ Về nơi bọn em chưa từng bị bỏ lại\dots''

Tiếng nói ấy tan trong không trung như hơi ấm cuối cùng của ngọn nến trước khi tắt phụt. Em ôm chặt con búp bê rơm vào lòng, như thể đó là mảnh ghép cuối cùng của một thời tuổi thơ bị đứt đoạn giữa chừng. Câu hỏi vang lên mang theo sự tủi thân tột cùng. Đó là khoảnh khắc của tôi - ``Kẻ đối diện''. Tôi bước tới bên cạnh Shino, nhìn thẳng vào đôi mắt xanh lục ấy.

Tôi bước sát lại, giọng trầm ấm: ``Anh chị không đủ mạnh để ngăn tai họa ấy xảy ra, và dĩ nhiên, không thể nào quay trở về quá khứ để thực hiện đựoc điều đó. Nhưng hôm nay, anh chị chắc chắn có thể đưa các em về nơi an nghỉ đích thực. Một nơi mà các em xứng đáng được tới.''

Tới đây, cô bé mắt xanh rưng rưng, giọt nước mắt trong suốt lăn dài trên má. Em ngẩng lên, đôi mắt xanh lục như chứa cả một đại dương đã cạn kiệt vì khóc.  Giọng em nhỏ đến mức chỉ còn là tiếng thì thầm của gió, nhưng mỗi chữ đều nặng trĩu thời gian:

``Các anh chị có biết không\dots\
Em đã từng đứng đây, đúng chỗ này, cách đây\dots\ rất lâu.
Khi ấy, trời cũng đêm như thế này, gió cũng rít qua kẽ lá như thế này.
Chỉ khác là\dots\ không có ai đứng về phía bọn em cả.''

Em dừng lại, nuốt nghẹn, tay siết chặt con búp bê rơm đến mức những sợi rơm rạ cũ kỹ bắt đầu đứt lìa.

``Họ\dots\ những người lớn ấy\dots\ bảo bọn em là `đứa trẻ mang tai họa'.
Họ bảo: `Nếu không hy sinh, cả làng sẽ chết đói'.
Em không hiểu `hy sinh' nghĩa là gì\dots\
Em chỉ thấy mẹ ôm em thật chặt, khóc đến nỗi vai áo bà ướt đẫm.
Rồi bà đẩy em ra, bảo: `Con phải mạnh mẽ, con là con của mẹ mà'.

\dots Rồi họ dẫn bọn em đi.
Có bạn khóc. Có bạn cắn vào tay người lớn, bỏ chạy.
Nhưng họ bắt lại, họ kéo bọn em bằng dây thừng\dots\ giống như khi nãy các anh chị thấy vậy.
Bọn em còn không có quyền phản kháng. Chỉ có thể nghe lời họ, tuân theo những gì mà người lớn bảo.
Em còn nhớ mùi đất bùn khi họ bắt bọn em quỳ xuống.
Em còn nhớ tiếng một bạn trai hét lên: `Mẹ ơi, con sợ lắm!'
Rồi tiếng ấy\dots\ tắt ngúm.

Em đã từng nghĩ: `Chắc mình sẽ gặp lại mẹ thôi.
Mẹ sẽ ôm mình, mẹ sẽ kể chuyện, mẹ sẽ bảo mọi thứ chỉ là ác mộng.'
Nhưng không.
Khi em mở mắt ra, chỉ có bóng tối.
Bóng tối\dots\ dày đến mức em có thể cảm nhận từng milimet nó đang bò vào da thịt mình.
Em không thể nhúc nhích.
Em không thể khóc, vì nước mắt cũng bị bóng tối uống cạn.

Em nghe thấy tiếng bạn bên cạnh thở\dots\ rồi thôi không thở nữa.
Em nghe thấy tiếng con búp bê của một bạn gái rơi xuống rồi im lặng.
Em nghe thấy\dots\ tiếng tim mình đập\dots\ cho đến khi nó cũng ngừng lại.''

Kuon khẽ gật gù, hai tay siết lại đến mức khớp xương trắng bệch; cậu cúi đầu, không dám nhìn thẳng vào đôi mắt xanh lục kia như thể chỉ cần một ánh nhìn cũng đủ khiến cậu cảm động: ``Bọn anh\dots\ thực sự xin lỗi,'' cậu đáp lại. ``Các em đã phải chịu đau khổ rồi.''

Shino đứng sát phía sau, hai vai rung lên nhè nhẹ; cô ôm chặt lấy ngực như thể muốn giữ lại trái tim đang vỡ vụn. ``Chị nghe thấy từng tiếng nấc của em,'' cô nức nở, nước mắt lăn dài trên má. ``Chị nghe thấy cả tiếng đất rơi\dots\ tiếng búp bê vỡ\dots\ và chị không thể làm gì. Chị xin lỗi\dots\ xin lỗi nhiều lắm.''

Sakura quỳ gối xuống đất, hai tay đặt lên đầu gối, mắt nhắm nghiền. ``Ranh giới này không chỉ để bảo vệ các em khỏi ma quỷ,'' cô thì thào, giọng nhỏ đến mức chỉ còn là hơi thở. ``Nó còn để bảo vệ các em khỏi chính sự tàn nhẫn của loài người\dots\ của chúng tôi. Xin hãy để chúng tôi chuộc lại lỗi lầm này, bằng cả đời này.''

Cô bé mắt xanh đó tiếp tục bộc lộ nỗi lòng của mình, bằng chất giọng trong trẻo như tiếng chuông gió, nhưng lại chứa đựng nỗi buồn của hàng ngàn năm:

``1300 năm\dots\
Các anh chị có hiểu 1300 năm là bao lâu không?
Là khoảng thời gian để cát bào mòn đá.
Là khoảng thời gian để một dòng sông đổi dòng.
Là khoảng thời gian để\dots\ một đứa trẻ không bao giờ lớn,
một lời xin lỗi không bao giờ được thốt,
một lời hứa không bao giờ được giữ.

Em đã từng đi lang thang trong bóng tối ấy,
gọi mẹ, gọi bạn, gọi trời, gọi đất\dots\
Chỉ để nghe lại tiếng vọng của chính mình:
`Các người đã bỏ rơi tôi.'

Nhưng mà, từ khoảnh khắc này, từ thời điểm các anh chị đứng trước mặt em,
em biết rằng\dots\ tất cả những điều đó sẽ thay đổi.
Khi các anh chị đứng đây,
khi các anh chị nói: `Chúng tôi nhớ các em.'
Em mới dám\dots\ khóc.
Em mới dám\dots\ tin.
Em mới dám\dots\ nói.

Cho nên là\dots\
nếu các anh chị có thể\dots\
đừng để bọn em\dots\
phải đợi thêm một 1300 năm nữa\dots\
được không?''

Phía sau chúng tôi, Nanase chắp tay cầu nguyện, Mari nhắm nghiền mắt không dám nhìn, còn Ayame cắn chặt môi, dường như đang phải kìm nén cảm xúc. Tới cả Yuko-sensei cũng không giấu được nét đượm buồn trên gương mặt sau khi nghe được những lời như vậy.

Sakura khẽ nghiêng đầu, nói khẽ như ru: ``Chúng tôi đã lập ranh giới này để không ai có thể làm tổn thương các em thêm lần nữa. Hãy tin chúng tôi.''

``Và nếu các em cần một bờ vai,'' Mari mở mắt, giọng run run, ``chúng tôi sẵn sàng chia sẻ.''

Cô bé mắt xanh chớp chậm đôi mi dài, như thể làn gió vừa thổi qua đã cuốn đi bụi bặm của mười ba thế kỷ. Đôi mắt lục bảo ấy trước chỉ phản chiếu bóng tối, nay bỗng loé lên một tia sáng mong manh, như ánh sao đầu tiên lọt qua khe nứt của hang động đóng băng.

Em không khóc ngay; trước hết là một tiếng nấc nghẹn lại trong cổ họng, rung nhẹ đầu vai bé xíu. Rồi, từng giọt nước mắt trong veo bắt đầu lăn, không ồ ạt, mà từng hạt, như thể thời gian đã dạy em phải tiết kiệm cả nước mắt, tất cả chỉ để cho ngày hôm nay, ngày mà linh hồn các em chính thức được siêu thoát.

Khi Kuon nói ``Các em không có lỗi'', em thả chậm con búp bê rơm xuống, ngón tay út vẫn siết lấy một sợi rạ tơi tả như níu kéo sự sống. Đầu em nghiêng nhẹ, mái tóc xanh đen rũ xuống che khuất đôi mắt đỏ ho, một cử chỉ của đứa trẻ lần đầu được phép e thẹn.

Azuki vừa dứt lời ``đưa về nhà'', em ngẩng lên, môi run run. Ánh trăng rọi thẳng vào gương mặt bé nhỏ: lúc ấy, nỗi đau không còn là vết thương mờ đỏ, mà biến thành một vì sao sáng rực trong đồng tử lục bảo.

Và rồi, em cười. Không phải nụ cười của kẻ được cứu rỗi, mà của đứa trẻ vừa nhận ra mình vẫn còn tên trong danh sách những người được yêu thương. Nụ cười ấy nhỏ đến mức chỉ có thể chứa được trong khoảnh khắc, nhưng đủ để cả khu rừng Aogami bỗng im lặng, như thể muôn loài đang cúi đầu trước một phép màu mong manh vừa được sinh ra giữa đêm dài.

``\textbf{Ranh giới đã mở, ký ức đã được thừa nhận!}'' Shion hét lên qua tiếng gió rít. ``Bây giờ chỉ cần cầu nối cuối cùng!''

Tôi lo lắng ngoái lại phía sau con đường mòn, khi người cuối cùng cho nghi thức này vẫn chưa xuất hiện. Nhanh lên nào, Suzu. Mọi người đang chờ em đấy\dots

\pov{Hikawa Suzuki}

Tôi cắm cúi chạy, hơi thở đứt quãng đau rát cả lồng ngực; mỗi nhịp nâng chân, mắt cá trái như có mũi tên nung đỏ đâm xuyên. Mồ hôi lạnh túa ra, trộn lẫn máu tươi từ vết xước trên má, nhỏ xuống cổ tay run rẩy. Phía sau, ba giọng hét xen lẫn nhịp thở dữ dội vang lên như trống trận:

``\textbf{Con nhãi kia! Đứng lại đó cho tao!}'' gã tóc vàng gầm vang, giọng khàn đặc vì tức giận và đau đớn.

``\textbf{Đại ca, nó đang chạy về hướng ánh đèn! Không được để nó tới kịp!}'' tên đàn em thứ nhất thở hồng hộc, tiếng bước chân giẫm lá khô rạo rực như muốn xé nát đêm.

Chúng lao vào nhau, va chạm cành cây, tiếng lá xào xạc như roi quất vào không khí. Mỗi lần tôi vấp gốc rễ, chúng hú lên đắc chí: ``Nó sắp ngã rồi! Nhanh lên!''

``\eng{\textbf{Chạy sang trái! Cúi đầu xuống!}}'' Game vang lên dứt khoát.

Tôi lập tức gập người. Một cành thông to bằng cổ tay quét sát đỉnh đầu, gió xoáy làm tóc dựng ngược. Phía sau, gã tóc vàng không kịp phản xạ, bị cành cây quật thẳng vào mặt. Tiếng *BỐP!* khô khan, hắn ngã lăn, máu mũi phụt ra, nhưng chỉ một giây sau đã gào: ``Mẹ kiếp con ranh!'' rồi vùng dậy, mắt đỏ ngầu, tiếp tục đuổi.

Hai tên đàn em càng thêm cuồng loạn. Chúng vừa chạy vừa chửi rủa từng họ hàng tôi, tiếng thở gấp gáp như chó săn bị thắt cổ. Mỗi khi có bụi cây cản, chúng đạp tung, lá bay vỡ tung như tiếng vỗ tay khích lệ đám quỷ.

Tôi tiếp tục bám theo ánh sáng xanh lấp lánh do các linh hồn trẻ thơ tạo ra. Dường như khu rừng đang cố tình mở đường cho tôi, trong khi liên tục dựng lên những chướng ngại vật đối với những kẻ truy đuổi. Một lớp sương mù đặc quánh bỗng tràn ra, cuộn lấy hai tên đàn em khiến chúng mất phương hướng, đâm sầm vào nhau và ngã lăn xuống một con dốc cạn.

Chúng lăn lông lốc, đầu đập đá, tiếng kêu thất thanh: ``Á! Đụng mẹ nó rồi!'' ``Mày đè lên chân tao!'' ``Đừng có buông! Bám vào!'' Nhưng chỉ một nhịp thở, chúng đã lồm cồm bò dậy, mắt hau háu, tiếp tục lao lên, máu me bê bết không kịp lau.

Tôi vừa nhảy qua một gốc cây bật gốc, tiếng bước chân của tên đàn em gần sát sau lưng. Bỗng đất dưới chân hắn nổi sóng, một dây rừng xanh biếc bật lên như con rắn thép, quấn cổ chân hắn lôi ngược trở lại. Hắn ngã chổng vó, mặt đập vào bùn. Game nói nhẹ trong đồng hồ: ``\eng{Root-level privilege acquired: trip-wire deployed.}''

Tên còn lại gầm lên xông tới, nhưng từng bước của hắn bỗng chìm sâu xuống đất, lá rừng mềm biến thành vũng lầy ảo. Từ trong lầy, bàn tay trẻ thơ phát sáng khẽ chạm vào ống chân hắn, thì thầm: ``Đừng theo\dots'' Tiếng thì thào vang dội như trống trận; hắn hoảng loạn lùi lại, vấp phải cây đổ mà Game đã bẻ gãy từ xa, ngã sõng soài. Đèn pin văng xa, tắt ngóm. Khu rừng khựng lại, nhường lối cho tôi băng qua trong im lặng, như mẹ hiền vỗ về đứa con gái nhỏ đang chạy trốn.

``\eng{\dots Vừa rồi là chơi chữ sao, Game-san?}'' tôi vừa chạy, vừa nhíu mày hỏi ông ấy.

``\eng{Dạng như vậy-- Mà, không còn thời gian hỏi đâu! Tiếp tục chạy về phía ánh đèn đi!}''

Tôi lao về phía trước, từng bước chân nện xuống đất như muốn đóng đinh mình vào con đường mòn. Gió rít qua tai, lá cây quất vào mặt tôi như những bạt tai lạnh lẽo. Phía sau, ba giọng hét vang lên đan xen, giòn giã như roi quất.

``\textbf{Đứng lại! Tao sẽ bẻ gãy chân mày!}'' gã tóc vàng gầm vang, tiếng thở gấp như bò rống.

Tôi cắn răng, bỏ qua cơn đau nhói ở cổ chân, lao qua một bụi cây thấp. Cành cây quật vào vai tôi, đau điếng, nhưng tôi không dừng. Ánh đèn pin phía trước lấp lánh như ngôi sao cuối chân trời, kéo tôi về phía nó bằng một sức mạnh vô hình.

``\textbf{Chạy nữa đi! Chạy nữa đi!}'' Game reo lên trong đồng hồ, giọng hưng phấn như đang cổ vũ một cuộc đua sinh tử.

Tôi gầm lên, dồn toàn bộ sức lực còn lại vào đôi chân. Mười mét. Tám mét. Năm mét. Tiếng hét phía sau càng lúc càng gần, nhưng ánh sáng trước mặt cũng rực rỡ hơn bao giờ hết.

Chẳng bao lâu sau, tôi vượt qua bờ rào sương mù cuối cùng, lao thẳng vào vòng sáng vàng ấm áp.

``Tới rồi!'' tôi reo lên khi nhìn thấy ánh sáng đèn pin le lói từ phía trước, cùng với bức tường năng lượng màu vàng của Sakura-san. Dù cổ chân đau như búa bổ, tôi vẫn bật thêm một nhịp chạy nước rút, mắt ướt đẫm vì xúc động và mừng rỡ: ``Chỉ cần thêm mười mét nữa\dots\ nghi lễ sẽ thành công!''

\pov{Hikawa Kaigou}

Từ trong làn sương mù mờ ảo, một bóng dáng nhỏ bé quen thuộc lao ra, quần áo xộc xệch và mái tóc rối bù.

``\textbf{Suzu!}'' tôi vội vã đưa tay ra. Kuon nhanh như cắt lao tới, nắm lấy tay con bé và kéo tuột vào bên trong ranh giới an toàn.

``Em xin lỗi vì đến trễ,'' Suzu thở dốc, gối khuỵu xuống nhưng ánh mắt hai màu của em vẫn rực sáng.

Chưa kịp để tôi hỏi han, từ phía con đường mòn, gã tóc vàng lảo đảo lao tới, theo sau là hai tên đàn em và cậu nam sinh rụt rè. Bọn chúng vằn mắt nhìn thấy Suzu, gầm lên tức tối: ``\textbf{Bắt lấy nó!}''

Nhưng bọn chúng đã chọn nhầm đối thủ.

``\textbf{Đừng hòng đụng vào đồng chí lớp tôi!}'' Ayame gầm lên. Hội trưởng Hội học sinh lao tới như một mũi tên, tung một cú đá xoay tuyệt đẹp trúng ngay giữa ngực gã tóc vàng, khiến hắn văng ngược trở lại ranh giới sương mù.

Gã tóc vàng mở to đôi mắt đỏ ngầu, kêu lên một tiếng ``\textbf{Ôi!}'' rồi đập lưng xuống đất, bụi lá bay tứ tung. Hắn ôm ngực, thở hổn hển, mặt nhăn nhó đau đớn, nhưng chỉ một giây sau đã gầm lên: ``Mẹ kiếp! Đánh tao à!'' và vội vã lồm cồm bò dậy, sẵn sàng lao vào trả đòn.

Hai tên đàn em thấy đại ca bị đá, liếc nhau một cái rồi hú lên như sói, xông thẳng vào Ayame. Tên thứ nhất giơ nắm đấm to tướng nện xuống đầu cô, nhưng Ayame khẽ lướt người, dùng đòn Aikido quật ngã hắn ra sau lưng. Tên còn lại rút từ túi quần ra một con dao gấp, múa may loạn xạ, lao tới như điên: ``Tao đâm chết cả lũ bây giờ!''

Nanase và Mari lập tức phối hợp. Mari ném mạnh chiếc đèn pin dự phòng vào mặt tên cầm dao, khiến hắn kêu ``Á!'' và loạng choạng che mặt. Nanase dùng chiếc đèn pin công suất lớn chiếu thẳng vào mắt hắn, ánh sáng trắng chói lòa làm hắn lảo đảo, dao rơi xuống đất.

Chứng kiến cảnh ba gã côn đồ bị phản kháng đầy quyết liệt, cậu nam sinh đi cùng thấy cảnh đó liền vội vàng lùi lại, giơ hai tay đầu hàng.

``Để bọn này lo vòng ngoài!'' Ayame lớn giọng, đứng chắn ngang con đường mòn với tư thế võ thuật vững chãi. ``Các cậu hoàn thành nghi thức đi!''

``\textbf{Làm đi, Suzuki!}'' Kuon hét lớn.

Suzu bước lên phía trước, đứng đối diện với miệng hố và cô bé mắt xanh. Con bé hít một hơi thật sâu, đôi mắt nhẹ nhàng lướt qua từng đốm sáng nhỏ đang lấp lánh quanh mình, giọng thì thầm đầy trìu mến:

``Các em không cô đơn nữa đâu\dots\ Cảm ơn các em đã chờ đợi. Giờ thì\dots\ chúng ta về nhà thôi.''

Rồi em giơ cao cánh tay đang đeo chiếc đồng hồ lên. ``\textbf{Game-san!}''

``Roger, madame. (Rõ rồi, thưa quý cô.)''

Chiếc đồng hồ phát ra một tiếng *BÍP* dài. Một luồng từ trường dịu nhẹ tỏa ra, kết nối bức tường năng lượng của Sakura, những hạt lúa của Shino và lời thừa nhận của tôi lại với nhau.

Sakura há hốc mồm, hai tay run run đặt lên ngực, mắt cô ánh lên niềm xúc động tột cùng: ``Kết nối\dots\ hoàn tất!''

Shino quỳ gục xuống đất, nước mắt lăn dài trên má, giọng thì thầm như cầu nguyện: ``Mong các em\dots\ hãy an nghỉ.''

Kuon siết chặt tay, mắt long lanh nhìn theo luồng sáng: ``Đây\dots\ là phép màu thật sự.''

Azuki ôm chặt lấy ngực, môi run run: ``Mình\dots\ mình không thể tin vào mắt mình nữa.''

Nanase và Mari đứng sững sờ, hai tay chắp lại trước ngực như đang chứng kiến một điều thiêng liêng.

Shion mở to mắt, giọng run run: ``Nghi thức\dots\ đang thành hình.''

Yuko-sensei đứng im lặng, hai mắt nhắm nghiền, dòng lệ lặng lẽ lăn xuống.

Gã tóc vàng mở to đôi mắt đỏ ngầu, môi run rẩy: ``Cá\dots\ cái quái gì thế này?''

Tên đàn em thứ nhất lùi lại hai bước, giọng lắp bắp: ``Chúng\dots\ chúng ta đang chứng kiến\dots\ ma thuật thật sao?''

Tên còn lại quỳ gục xuống đất, hai tay ôm đầu, thét lên: ``Không! Không thể nào! Nó không thể xảy ra ở ngoài đời thực được!!!''

Cậu nam sinh rụt rè đứng chết lặng, mắt tròn xoe đầy kinh ngạc và gửi gắm một chút hy vọng.

Suzu khép hờ đôi mắt, nói bằng tất cả sự chân thành: ``Đường đã mở rồi. Các em không còn bị trói buộc ở nơi này nữa. Đi thôi, về nhà thôi.''

Lời nói vừa dứt, cô bé mắt xanh mỉm cười. Nụ cười tươi sáng và hồn nhiên nhất của em. Em từ từ buông con búp bê rơm xuống lòng hố, cơ thể em bắt đầu phát sáng và tan biến thành những hạt bụi lấp lánh. Từ dưới lòng hố, hàng ngàn đốm sáng nhỏ li ti bay vút lên bầu trời đêm, tạo thành một dải ngân hà rực rỡ chảy ngược lên không trung.

Khắp không gian vang lên những tiếng cười khúc khích trong trẻo, không còn sự oán hận, chỉ còn lại niềm vui của những linh hồn được giải thoát, được tề tựu nơi thiên đàng. Những đốm sáng bay lượn quanh chúng tôi như một lời cảm ơn, trước khi tan biến hoàn toàn vào màn đêm tĩnh lặng.

Bức tường năng lượng của Sakura hạ xuống. Sương mù tan biến, để lộ ra một khu rừng bình yên dưới ánh trăng. Ba gã côn đồ chứng kiến toàn bộ cảnh tượng thần thánh đó, sợ hãi đến mức ngất xỉu ngay tại chỗ.

Shino nhìn theo những đốm sáng cuối cùng bay đi, khóe môi khẽ nhếch lên một nụ cười mãn nguyện. ``Cuối cùng\dots\ mọi chuyện cũng đã kết thúc.''

\pov{Hikawa Kuon}

Tôi thở phào một hơi dài, cảm giác như một tảng đá ngàn cân vừa được nhấc khỏi lồng ngực. Nhìn những đốm sáng cuối cùng tan biến, khu rừng Aogami bỗng trở nên tĩnh lặng một cách lạ thường, không còn sự nặng nề hay u ám của những oán linh kéo dài hàng thế kỷ.

Nhưng rồi, sự nhẹ nhõm đó không kéo dài lâu. Không, không phải vì những tinh linh bé nhỏ kia. Mà là vì một cái gì đó khác từ chúng tôi.

``Suzu! Em ổn chứ?'' Kaigou lao vụt tới, quỳ thụp xuống bên cạnh em gái mình, người vẫn đang thở dốc vì kiệt sức.

``Em\dots\ em ổn, Onee-sama,'' Suzuki mỉm cười yếu ớt, cố gắng đứng dậy nhưng mắt cá chân sưng tấy khiến em loạng choạng.

Cùng lúc đó, ba tên côn đồ nằm bẹp dưới đất cũng bắt đầu lầm bầm tỉnh lại. Ánh sáng từ dải ngân hà linh hồn thăng thiên dường như vẫn còn đọng lại trong võng mạc của chúng, khiến chúng ngơ ngác như những kẻ vừa từ cõi chết trở về.

``Thiên\dots\ thiên đường sao?'' tên cầm đầu tóc vàng thều thào, ánh mắt vẫn còn thất thần. ``Tôi vừa thấy\dots\ những thiên sứ\dots''

Tôi khẽ nhếch mép. Thiên sứ cái gì chứ? Nếu chúng biết chúng suýt nữa đã bị những ``thiên sứ'' đó lôi xuống huyệt mộ thì chắc không dám thốt ra câu đó đâu.

Trong khi đó, Kaigou đang lo lắng kiểm tra tay chân cho Suzuki - cũng không trách được, em ấy là người mà cô ấy hết mực yêu thương. Khi em ấy khẽ nhăn mặt đau đớn lúc Kaigou chạm vào vai, cô ấy lật nhẹ lớp áo khoác mỏng của em lên. Một vết bầm tím lớn hiện rõ trên bả vai trắng ngần của Suzuki.

``\dots Ôi trời.''

Trong tích tắc, bầu không khí xung quanh Kaigou đột ngột giảm xuống độ không tuyệt đối. Sát khí tỏa ra từ cô nàng đặc quánh đến mức tôi cảm thấy gai người, dù tôi chẳng làm gì sai. Cô ấy không nói lời nào, nhưng đôi mắt màu xanh nước biển ấy giờ đây tối sầm lại, sâu thẳm và đáng sợ như một vực thẳm không đáy.

Tôi nhìn sang ba tên côn đồ đang lồm cồm bò dậy, khẽ lắc đầu thở dài: ``Các cậu vừa làm gì với Suzuki khiến em ấy trông thảm hại như vậy?''

Để đổ thêm dầu vào lửa, Game từ chiếc đồng hồ trên tay Suzuki đột ngột lên tiếng. Giọng nói vô cơ của ông ta vang vọng giữa khu rừng đêm: ``Nếu cô muốn biết, Kaigou-user, thì những kẻ hạ đẳng này đã trói em gái cô lại, định quay clip làm nhục và thậm chí là rạch nát mặt cô bé. Nếu không có `rừng phản ứng' và sự can thiệp của tôi, có lẽ giờ này Suzuki-user đã không còn lành lặn để đứng đây đâu.''

Nghe đến đó, không chỉ Kaigou mà cả nhóm Nanase, Ayame và Yuko-sensei đều sững sờ, rồi ngay lập tức chuyển sang trạng thái phẫn nộ tột cùng.

``\dots Thôi, xin vĩnh biệt các cậu. Tội tày đình rồi,'' tôi lắc đầu ngán ngẩm, giọng nói bình thản nhưng đầy ẩn ý. ``Các cậu đã tự đào mồ chôn mình rồi đấy.''

``Tôi không ngờ các cậu lại có thể đê hèn đến mức đó đấy,'' Nanase nghiến răng, đôi mắt sắc lẹm nhìn ba tên kia.

Yuko-sensei đặt nhẹ tay lên vai Kaigou, ánh mắt buồn bã nhưng lạnh lùng: ``Là giáo viên, tôi đã thất bại khi để học sinh mình gặp nguy hiểm. Nhưng từ giây phút này\dots\ tôi sẽ không cho phép bất cứ ai động đến các em thêm một lần nào nữa.''

Kaigou hít một hơi thật sâu, giọng cô nhỏ nhẹ nhưng mang theo uy lực lạnh thấu xương: ``Ayame-san, Nanase-san. Với tư cách là Hội trưởng Hội học sinh và tiểu thư nhà Furukawa đang có mặt tại hiện trường\dots\ tôi có quyền đánh túi bụi bọn chúng không?''

Ayame khoanh tay, khuôn mặt vốn nghiêm nghị giờ đây càng thêm đanh lại. Cô nàng nhìn sang Nanase, rồi nhìn sang ba tên côn đồ đang run cầm cập trước áp lực của cô chị lớn.

``Mặc dù theo quy định, trường không cho phép học sinh dùng bạo lực,'' cô nàng tóc trắng cất giọng đều đều, ``nhưng với tư cách Hội trưởng, tôi xác nhận hành vi của các cậu là cực kỳ nghiêm trọng, đe dọa đến tính mạng và danh dự học sinh khác. Tôi sẽ\dots\ tạm thời quay đi hướng khác trong năm phút tới.''

Nanase gật đầu, môi nhếch lên một nụ cười lạnh: ``Nhà Furukawa cũng cam kết sẽ không nhìn thấy gì cả. Sakura-san, Shino-san, các cậu thì sao?''

``Tôi không thấy gì hết,'' Sakura quay mặt đi, lẩm nhẩm về việc kiểm tra ranh giới. Shino cũng gật đầu im lặng.

Yuko-sensei khoanh tay, thở dài một hơi rồi nhìn lên tán cây: ``Đêm nay rừng tối quá, tôi chẳng nhìn rõ ai đang làm gì ai đâu. Các em thích làm gì thì làm.''

Shion thì bồi thêm bằng giọng chuunibyou đặc trưng: ``Các ngươi có biết hành động trời không dung đất không tha kia suýt chút nữa đã khiến nghi lễ đại bại không!? Nếu không có người kết nối, cả năm người chúng ta đã bị oán khí nuốt chửng rồi. Đó gọi là tội ngộ sát đấy, lũ phàm phu tục tử!''

``Cậu nói đúng đó, Shion-san,'' Nanase tiếp lời. ``Tội ngộ sát và bắt cóc trẻ vị thành niên đủ để các cậu bóc lịch vài năm đấy.''

Ba tên côn đồ giờ mới thực sự tỉnh táo, mặt cắt không còn giọt máu. Tên tóc vàng lắp bắp: ``Chúng\dots\ chúng tôi xin lỗi! Tại chúng tôi chỉ định dằn mặt thằng Kuon thôi!''

``Dằn mặt tôi?'' tôi nhướng mày. ``Vậy thì các cậu nên nhắm vào tôi, thay vì một cô gái không có khả năng tự vệ. À mà, có được đâu. Các cậu còn chẳng đấm nổi cái nào vào người tôi nữa là.'' Câu sau tôi nói có phần mỉa mai hơn.

Ngay trước khi cô nàng mắt xanh nước biển bắt đầu thi hành án, lúc này, Suzuki khẽ lên tiếng, chỉ tay về phía cậu nam sinh rụt rè đang đứng khúm núm một góc: ``Đợi đã\dots\ anh bạn kia không làm gì em cả. Cậu ấy bị ép phải đi theo và đã bí mật cởi trói cho em để tôi chạy thoát đấy ạ.''

Kaigou quay lại, giọng trầm ấm nhưng đầy uy lực: ``Cậu kia, ra đây.'' Cậu nam sinh run rẩy bước tới, hai tay giấu sau lưng, đầu cúi sát ngực. ``Có thật không?''

``Tôi\dots\ tôi xin lỗi,'' cậu thì thào, giọng nghẹn lại. ``Họ bắt tôi phải đi cùng\dots\ nhưng khi thấy Suzuki-san khóc, tôi\dots\ tôi không chịu nổi. Tôi đã lén cắt dây trói cho cậu ấy\dots\ Tôi biết mình hèn, nhưng\dots\ tôi thực sự không muốn làm tổn thương ai cả.''

``Xác nhận,'' Game nói thêm, ``cậu nhóc này là người duy nhất còn sót lại chút lương tri trong đám rác rưởi kia.''

Nghe những gì mà ba người trình bày, Ayame gật đầu đồng tình: ``Vậy cậu có thể sang đứng cạnh Yuko-sensei. Còn ba người các cậu\dots''

Kaigou bước tới trước mặt ba tên côn đồ, đôi bàn tay của cô nàng bẻ khớp kêu răng rắc khiến cho bọn chúng giật thót. Ánh trăng chiếu rọi vào khuôn mặt lạnh lùng của cô, tạo nên một vẻ đẹp đầy nguy hiểm. ``Chúng mày nghe cho rõ đây,'' cô gằn giọng, xưng hô hoàn toàn thay đổi. ``Chúng mày có thể đụng vào Kuon-kun, tao có thể xem xét bỏ qua. Đụng vào tao, may ra chúng mày cũng chỉ bay mất mấy cái răng. Nhưng động vào em gái tao\dots\ là sai lầm lớn nhất và cuối cùng trong cuộc đời học sinh của chúng mày đấy.''

Tên tóc vàng run rẩy, môi mấp máy: ``Xin lỗi\dots\ xin lỗi nhiều lắm\dots\ tôi thực sự hối hận\dots\ Làm ơn đừng đánh tôi!!''

Cô nàng lạnh lùng cắt ngang: ``Tao đếch cần cái lời xin lỗi tởm lợm ấy của mày. Và yên tâm, tao mà đánh thì nhẹ tay lắm. Sương sương khoảng\dots\ trăm cú thôi.''

Nhìn thấy sắc mặt biến dạng của cả ba kẻ côn đồ xui xẻo, Game đột ngột phát ra tiếng cười khẩy: ``\eng{Any last words?}\ (Có lời trăn trối nào không?)''

``{\Large\textbf{CỨU MẠNG!!!}}''

Tiếng hét của ba tên côn đồ vang vọng khắp khu rừng Aogami, xé toạc màn đêm tĩnh lặng. Nhưng đúng như kịch bản, chẳng có một ai ở đây mảy may mủi lòng. Những cú đấm, cú đá của Kaigou trút xuống như mưa. Sức mạnh của cô ấy mạnh thế nào, có lẽ khỏi cần phải bàn, và khi cô nàng điên tiết, đó thực sự là một thảm họa. Tôi biết việc đánh người như này là điều cực kỳ không khuyến khích, nhưng tôi biết rõ cô ấy thương yêu Suzuki tới mức nào; và với một người chị, việc ra tay bảo vệ em gái mình khỏi những hiểm họa là điều mà tôi có thể thông cảm.

Cơ mà, có lẽ sau khi trở về chắc tôi phải huấn luyện lại Kaigou vài buổi để kiềm chế cảm xúc hơn. Như bố tôi từng nói: ``Nếu con mà cứ cáu gắt như thế này thì ra đời không làm được gì đâu con ạ.''

Tôi đứng ngoài nhìn cảnh tượng đó, rồi quay sang hỏi Ayame: ``Này Hội trưởng, tôi cứ ngỡ cậu chỉ giỏi sổ sách với chơi khăm mọi người thôi, không ngờ cậu cũng biết võ thuật và ra đòn dứt khoát vậy à?''

Cô nàng đáp lại với giọng tươi tỉnh như mọi khi dù tiếng la hét phía sau vẫn đang tiếp diễn: ``Nhà tôi vốn có truyền thống võ đạo. Tôi đã học aikido từ năm sáu tuổi để tự vệ và duy trì kỷ luật bản thân. Đôi khi bạo lực là cần thiết để trị mấy kẻ cứng đầu như thế đó\~{}''

``Ra là vậy,'' tôi lẩm bẩm. May mà mình chưa bao giờ làm gì đắc tội với cô nàng này.

Tôi khẽ kéo tay áo Shion, kéo cô nàng ra xa đám đông một chút rồi thì thầm: ``Này, hỏi chút. Nếu nghi thức vừa rồi thất bại, bọn tôi có thật sự\dots\ chết không?''

Cô nàng chớp mắt, rồi bật cười khúc khích như vừa nhớ ra chuyện gì đó tròn mắt: ``Ế? À, không đâu. Ta chỉ nói vậy cho nó căng thôi.'' Cô nàng lấy trong túi ra một cuốn sổ nhỏ đã cũ, ghi gì đó rồi đưa cho tôi xem. ``Tài liệu có ghi rằng, nghi thức này chỉ thực hành một lần mỗi năm thôi. Nếu lỡ trượt, năm sau làm lại được mà.''

Tôi trợn mắt, vừa thở phào vừa mếu máo: ``Vậy tức là\dots\ tôi có thể được thực hiện lại sao?''

``Ừ, nhưng mà\dots'' Shion gãi má, nhìn lên trăng, ``nếu năm sau vẫn thất bại, thì cứ thế lặp lại thôi. Cứ mỗi năm một lần, cho đến khi thành công.''

Tôi ôm trán, thì thầm như than trời: ``Không biết tôi có đủ sức chờ thêm một năm nữa không đây\dots\ Cậu mà nói sớm thì tôi đỡ lo đến mức suýt tái mét cả mặt.''

Nàng Phù thủy bỗng nghiêng đầu, mắt sáng quắc: ``Kuon, cái giọng vừa cất lên trong đồng hồ ngươi là ai vậy? Nghe như\dots\ linh hồn bị phong ấn trong cổ vật ngàn năm!''

Kuon thở dài, gãi má: ``À\dots\ để giấu cũng mệt. Thôi, xin giới thiệu luôn, đây là Game, trợ lý ảo của tôi.''

Ngay lập tức, chiếc đồng hồ phát ra tiếng *TÁCH* nhỏ, giọng ông hiện rõ sự mệt mỏi vang lên giữa rừng đêm: ``Yeah, yeah. Tôi là Game, AI hỗ trợ đa nhiệm. Phải, tôi biết tên tất cả mọi người: Shitou Shion, Furukawa Nanase, Onidori Ayame, Akagane Mari, Utachi Azuki, Tsukishita Yuko, và cả Shinomiya Sakura, Misora Shino nữa. Rất hân hạnh được gặp.''

Shion giật mình, hai tay che miệng, mắt long lanh: ``Oa! Mọi người nghe chưa? Hắn ta gọi đúng danh xưng thật của ta - Shion, nữ pháp sư bóng tối!'' Cô quay sang Kuon, hạ giọng: ``Bộ\dots\ bộ ngươi đang nuôi một con quỷ công nghệ trong đồng hồ sao?''

Nanase cau mày, đẩy nhẹ gọng kính: ``Trợ lý ảo mà biết cả họ tên học sinh? Kuon-san, cậu có thể giải thích rõ hơn không?''

Kuon gãi đầu: ``Ờ\dots\ thì nó quét dữ liệu công khai của trường, rồi ghép với camera trên đồng hồ. Đừng lo, ông ấy\dots\ tương đối vô hại.''

Nàng Phù thủy tự xưng thở dài, nhún vai: ``Từ lúc ngươi độc thoại một mình trong lớp là ta thấy nghi nghi rồi. Hóa ra là vì cái thực thể trong đó.''

Sakura cười trừ: ``Tôi cũng từng bắt gặp ông ấy rồi.'' Cô khẽ gật về phía đồng hồ: ``Chào ông nha, Game-san.''

Game đáp lại bằng giọng đầy ẩn ý: ``Luôn sẵn sàng hỗ trợ, Sakura-san. Và xin lưu ý, lũ côn đồ kia đang tỉnh dậy, nên giữ khoảng cách an toàn.''

Shion reo lên, giơ cao tay: ``Tuyệt! Từ nay ta có thêm đồng minh công nghệ! Game, ngươi có thể phát hiện dấu hiệu ác quỷ không?''

``Tùy mức độ ác quỷ cô định nghĩa,'' Game đáp, ``nhưng nếu là kẻ định quay clip làm nhục học sinh để thỏa mãn cái tôi cao ngạo của kẻ bắt nạt, thì\dots\ 100\%.''

Nanase thở dài, khoanh tay: ``Được rồi, Kuon-san, sau này cậu sẽ phải ký cam kết bảo mật với Hội học sinh. Còn bây giờ, mọi người, tập trung, chúng ta còn việc phải làm.''

Chúng tôi tiếp tục hướng về cô nàng tóc đuôi ngựa đang tiếp tục quất những đòn roi vun vút lên những kẻ côn đồ đang van xin tha mạng kia. Không biết cô nàng đã thấm mệt chưa nhỉ?

Azuki đứng cạnh tôi, khẽ tặc lưỡi khi thấy gã tóc vàng bị Kaigou quật ngã lần thứ bao nhiêu không đếm xuể: ``Chà, Kaigou-chan đúng là không nương tay chút nào. Nhìn cậu ấy đánh kìa, chắc mai bọn kia phải nhập viện mất.''

Mari thì vừa sợ vừa hả dạ: ``Đáng đời lắm! Cho chừa cái tội bắt nạt Suzuki-san.''

``Công lý đã được thực thi,'' Game vang lên đầy hài hước độc địa, tiếng tí tách như đang gõ bàn phím ảo. ``Có cần tôi tăng sức mạnh cho cô để cho cô đánh thả ga hơn không, Kaigou-user?''

Yuko-sensei nhìn ba tên côn đồ nằm co ro, khẽ thở dài, giọng trầm ấm nhưng đầy nghiêm khắc: ``Không cần đâu, Game-san. Đôi khi\dots\ bài học đau đớn nhất lại là liều thuốc hữu hiệu nhất.'' Cô đưa mắt lên trăng, lẩm nhẩm: ``Hy vọng lần này\dots\ chúng sẽ nhớ đời.''

Sau khoảng mười phút ``xử lý'' không nương tay, Kaigou mới chịu dừng lại khi hơi thở đã bắt đầu dồn dập. Ba tên côn đồ nằm la liệt trên đất, mặt mũi sưng vù, quần áo rách tướp và khóc không thành tiếng. Oái oăm thay, chính Suzuki lại là người bước tới, đưa cho chúng mấy miếng băng cá nhân từ túi sơ cứu của mình. Một sự sỉ nhục ngọt ngào hơn bất cứ đòn đánh nào.

``Tôi hy vọng các anh sẽ ghi nhớ bài học này,'' Suzuki nói nhỏ nhẹ, nhưng ánh mắt em giờ đây đã không còn sự rụt rè trước đây nữa. Cả ba tên côn đồ đó hầu như không cất lên nổi giọng nói nào sau trận đòn gần như vô tận này.

\ct

Ánh trăng bàng bạc xuyên qua kẽ lá, rải những vệt sáng lung linh lên con đường mòn vốn trước đó chỉ toàn bóng tối và nỗi sợ. Chúng tôi bắt đầu hành trình rời khỏi cánh rừng Aogami. Tôi khom người, xốc Suzuki lên lưng mình một cách cẩn thận. Mắt cá chân của em đã sưng vù như một trái mận chín, và mỗi nhịp thở của em phả vào gáy tôi đều mang theo sự mệt mỏi cùng cực.

``Em có nặng không, Onii-sama?'' Suzuki thầm thì, giọng em nhỏ xíu như tiếng gió thoảng.

Tôi mỉm cười, bước đi chậm rãi nhưng chắc chắn trên nền lá mục: ``Em nhẹ hều hà, Suzuki. Có khi còn nhẹ hơn cả cái ba lô đầy sách của anh đấy.'' Tôi thầm nghĩ, vết thương của em may mắn không quá nghiêm trọng, nhưng bác sĩ chắc chắn sẽ bắt em phải nghỉ ngơi vài ngày mới đi lại bình thường được.

Bên cạnh tôi, Kaigou đang sải bước với một vẻ sảng khoái chưa từng thấy. Cô ấy hít một hơi thật sâu không khí trong lành của rừng đêm, đôi tay thỉnh thoảng lại khẽ cử động như vẫn còn cảm nhận được dư chấn của trận ``vận động'' vừa rồi.

``Đánh xong thấy người nhẹ nhõm hẳn luôn ấy mọi người ạ\~{}'' Kaigou cười rạng rỡ, quay sang nhìn nhóm bạn. ``Cứ như là trút được hết đống bực dọc tích tụ cả tháng trời qua vậy.''

``Cậu thì nhẹ nhõm rồi, chỉ khổ cho ba tên kia thôi,'' Azuki vừa đi vừa tặc lưỡi, tay vẫn còn cầm chiếc đèn pin chiếu đường. ``Nhưng mà công nhận, Kaigou-chan ra đòn chuẩn thật đấy. Mình đứng xem mà còn thấy đau giùm.''

Mari gật đầu lia lịa, vẻ mặt vẫn còn chút phấn khích: ``Đúng đúng! Nhìn gã tóc vàng kia ngã chổng vó mà tôi thấy hả dạ cực kỳ. Suzuki-san của chúng ta không phải hạng người để các anh muốn bắt nạt là bắt nạt đâu!''

Yuko-sensei đi phía sau cùng, ánh mắt cô dịu lại khi nhìn thấy sự đoàn kết của cả nhóm. ``Dù sao thì bài học này cũng đủ để chúng nhớ đời rồi. Các em đã làm rất tốt, không chỉ là cứu Suzuki-san mà còn hoàn thành nghi lễ một cách xuất sắc.''

``Nghi lễ đã thành công, ranh giới đã được phong ấn trở lại,'' Sakura khẽ vuốt cành cây ven đường, nụ cười nhẹ nhàng nở trên môi. ``Khu rừng giờ đây không còn tiếng khóc oán than nữa. Cảm ơn mọi người đã tin tưởng tôi.''

Shino đi bên cạnh Sakura, ánh mắt cô cũng trở nên bình yên hơn bao giờ hết: ``Ký ức của chúng đã được xoa dịu. Mình có thể cảm nhận được hơi thở của rừng đã thay đổi. Nó không còn lạnh lẽo nữa, mà trở nên\dots\ bao dung hơn.''

``Hừm, sức mạnh bóng tối của ta kết hợp với linh hồn công nghệ của Game đúng là vô đối!'' Shion giơ cao tay, áo choàng (thực ra là áo khoác) bay phấp phới dù trời chẳng mấy gió. ``Game, ngươi thấy sao về màn trình diễn của 'Nữ pháp sư bóng tối' này?''

Chiếc đồng hồ trên tay Suzuki phát ra tiếng cười khẩy: ``Nếu cô định nghĩa việc đứng hét lớn và ném bùa là 'màn trình diễn', thì tôi cho 10 điểm về thần thái, Shion-san. Nhưng về hiệu quả thực tế, có lẽ tôi nên cộng điểm cho cú đá của Ayame-san nhiều hơn.''

Ayame cười tít, vẻ mặt vẫn giữ được sự nhây thường thấy: ``Hơm dám, hơm dám\~{} Tôi chỉ làm những gì cần thiết để bảo vệ học sinh của mình thôi. Mà, Kuon-kun này, sau khi trở về, tôi cần cậu nộp một bản báo cáo chi tiết về sự kiện này. Nanase-san cũng sẽ hỗ trợ việc lưu trữ hồ sơ bảo mật để tránh rò rỉ thông tin tâm linh ra ngoài.''

Nanase gật đầu, ánh mắt cô lướt qua khu rừng một lần cuối: ``Phía gia đình Furukawa cũng sẽ cử người đến kiểm tra định kỳ để đảm bảo không có sự cố nào phát sinh. Chúng ta đã làm xong phần việc của mình. Giờ là lúc để Aogami tự chữa lành.''

Chúng tôi cứ thế vừa đi vừa trò chuyện, tiếng cười nói vang vọng giữa đại ngàn yên tĩnh. Khu rừng Aogami đêm nay không còn là nơi của nỗi sợ, mà là nhân chứng cho tình bạn, tình thân và sự dũng cảm của chúng tôi. Mỗi bước chân rời đi, tôi cảm nhận được sự thanh thản kỳ lạ lan tỏa khắp cơ thể. Mọi chuyện, thực sự đã ổn rồi.

\begin{center}
  \uline{Ngày hôm sau.}
\end{center}

Không khí ở trường trung học Aogami đặc quánh sự tò mò và những lời đồn đại. Suzuki đã trở lại trường, nhưng trong một tình trạng khiến ai đi ngang qua cũng phải ngoái nhìn: chân trái của em được nẹp cứng và em phải di chuyển bằng hai chiếc gậy chống. Em ấy không tỏ vẻ khó chịu gì, nhưng chị gái của em thì không nghĩ như vậy. Thi thoảng, Kaigou còn xua đuổi những người hiếu kỳ sán tới, kèm với câu kiểu: ``Bộ các người chưa từng thấy ai bị gãy chân bao giờ à?'' khiến tôi và Suzuki không kìm được mà bật cười khúc khích.

Sáng sớm hôm đó, Kaigou đã càm ràm không ngớt, bắt em phải ở nhà nghỉ ngơi. Cô ấy thậm chí còn định khóa cửa nhốt em lại, nhưng Suzuki đã một mực lắc đầu: ``Em muốn thấy kết cục của chuyện này, Onee-sama. Với lại, em không muốn vì mấy vết thương này mà bỏ lỡ bài vở trên lớp đâu.'' Sự cứng đầu của em đôi khi còn vượt xa cả tôi và Kaigou cộng lại. Nếu Suzuki là bướng bỉnh số hai, chắc chỉ có mẹ tôi mới là đối thủ xứng tầm với em ấy.

Trong buổi tập trung chào cờ đầu tuần, cả sân trường chìm trong sự im lặng lạ thường khi Ayame-san bước lên bục giảng với một xấp hồ sơ dày cộm. Trước mặt toàn thể các bạn học sinh, cô bắt đầu đọc to những tội trạng mà ba kẻ côn đồ kia đã gây ra. Cứ mỗi tội được đọc, bên dưới các học sinh lại có những tiếng bàn tán xôn xao. Cuối cùng, tới bản quyết định.

``Căn cứ vào những bằng chứng xác thực và lời khai của các nhân chứng,'' giọng Ayame vang vọng qua hệ thống loa, đanh thép và không chút khoan nhượng, ``Hội học sinh và Ban giám hiệu nhà trường quyết định đình chỉ học vô thời hạn đối với ba học sinh lớp 3-D do hành vi bắt cóc, đe dọa hành hung và xâm phạm danh dự học sinh lớp dưới. Hồ sơ vụ việc hiện đã được chuyển giao cho cơ quan cảnh sát địa phương để tiếp tục điều tra.''

Cả sân trường ồ lên một tiếng đầy kinh ngạc, rồi ngay lập tức chuyển thành những tiếng xì xào bàn tán dữ dội. Ba tên côn đồ bị gọi mặt chỉ tên phải đứng ra trình diện trước toàn trường, đầu cúi gầm đến mức cằm chạm vào ngực. Những vết bầm tím trên mặt chúng vẫn còn rõ mồn một - dấu ấn từ ``bài học'' của Kaigou đêm qua. Chúng không dám ngẩng mặt nhìn ai, đôi vai run rẩy trong sự nhục nhã tột cùng trước hàng nghìn con mắt đang đổ dồn vào mình.

Tôi đứng ở hàng lớp 1-C, khoanh tay thở phào nhẹ nhõm. Quay sang bên cạnh, Kaigou đang nhìn về phía ba tên kia với ánh mắt sắc lẹm, đôi môi khẽ nhếch lên một nụ cười lạnh lùng hệt như một nữ vương vừa dẹp xong loạn lạc. Còn Suzuki, em đứng tựa vào đôi gậy chống, khẽ thở phào và gật đầu nhẹ như một lời tiễn biệt cho quá khứ không vui vẻ.

Đó là lần cuối cùng bộ ba Hikawa chúng tôi nhìn thấy ba khuôn mặt đó trong trường. Chúng biến mất như chưa từng tồn tại, để lại một bài học đắt giá cho bất cứ kẻ nào dám có ý định đụng vào gia đình chúng tôi.

Tôi đứng dưới sân trường, nhìn lên bầu trời xanh ngắt, khẽ mỉm cười. Mọi chuyện, cuối cùng cũng đã được giải quyết êm đẹp.

\ct

Những ngày tiếp theo sau nghi lễ tại rừng Aogami, thị trấn dường như khoác lên mình một lớp áo mới, rạng rỡ và chân thực hơn bao giờ hết. Nếu ở những ngày đầu đặt chân tới đây, Aogami hiện lên với vẻ thanh bình nhưng luôn ẩn chứa một cái gì đó nằng nặng, bí ẩn len lỏi trong mùi gỗ ẩm và những tán cây rừng xanh thẳm, hay như lần hồi sinh sau vòng lặp cũ mang theo sự ngỡ ngàng của một giấc mơ vừa tỉnh; thì giờ đây, sự thay đổi mang tính bản chất.

Không khí không còn cái cảm giác lạnh lẽo, gai người đặc trưng mỗi khi hoàng hôn buông xuống - thứ oán khí vô hình mà tôi từng phải dùng công việc bán thời gian để quan sát. Những bóng đen mờ ảo vốn thường lảng vảng nơi góc phố, những tiếng thì thầm vô định trong gió hay cảm giác bị theo dõi từ những bụi rậm dọc con đường lát đá uốn lượn đều đã hoàn toàn biến mất. Thay vào đó là một thực tại sáng sủa và trong trẻo tuyệt đối.

Ánh nắng buổi sớm rọi xuống những mái nhà gỗ truyền thống nằm xen kẽ với các tòa nhà bê tông hiện đại, tạo nên một bảng màu rực rỡ và ấm áp. Mùi nhựa thông phảng phất trong gió xuân không còn lẫn vị mục ruỗng của thời gian bị giam cầm. Tiếng xe điện chạy êm ru trên những con phố dưới chân núi giờ đây nghe thật vui tai, hòa cùng tiếng người gọi nhau mở hàng buổi sớm đầy sức sống. Aogami hiện tại không còn là một ``bình chứa'' cho những linh hồn bị bỏ lại, mà đã thực sự trở thành một ngôi nhà bình yên cho những người đang sống. Những cánh hoa hồng phấn bay trong gió không còn mang theo sự ma mị của một kỳ tích gượng ép, mà là sự rạng rỡ của một khởi đầu mới thực sự. Có thể nói, thị trấn Aogami đã hoàn toàn được tái sinh, cả về nghĩa vật lý lẫn tâm linh.

Cuộc sống học đường của chúng tôi trở lại quỹ đạo bình thường. Những tiết học, những giờ giải lao ồn ào và cả những buổi cuối tuần rộn rã. Căn nhà của bộ ba Hikawa giờ đây đã trở thành ``đại bản doanh'' của lớp 1-C chúng tôi. Mari, Shion, Ayame và thỉnh thoảng là cả Nanase hay Sakura đều tụ tập ở đây để cùng học bài, nấu ăn hoặc đơn giản là nghe Shion luyên thuyên về những ``truyền thuyết bóng tối'' mới mà cô ấy vừa nghĩ ra. Những bạn học khác như Inari, Miku, Akane và Yuka thi thoảng kéo tới làm vài ván trò chơi điện tử, hoặc những người như Honoka và Suzune rủ tôi làm những trò ``con bò.''

Vết thương ở chân của Suzuki cũng dần lành lại. Em đã có thể bỏ gậy và đi lại bình thường, dù Kaigou vẫn còn hơi quá bảo bọc mỗi khi thấy em bước đi hơi nhanh. Về phần mình, tôi đã chính thức xin nghỉ công việc bán thời gian tại cửa hàng tiện lợi. Mục đích ban đầu của tôi khi làm ở đó là để thu thập thông tin và quan sát những thực thể oán linh ẩn mình trong bóng tối thị trấn. Giờ đây, khi Aogami đã yên bình, lý do đó không còn tồn tại nữa. Tôi cũng đã dành dụm được một số tiền kha khá cho công việc này, và giờ cũng chưa nghĩ được sẽ tiêu cho việc gì.

Việc hậu sự cũng tốn không ít thời gian. Tôi đã phải thức trắng vài đêm để hoàn thành bản báo cáo chi tiết dài mười trang cho Hội học sinh theo yêu cầu của Ayame. Nanase thì làm việc cực kỳ chuyên nghiệp, cô ấy đã mã hóa toàn bộ dữ liệu về nghi lễ và lưu trữ chúng vào hệ thống bảo mật của gia đình Furukawa, đảm bảo rằng bí mật về oán linh rừng Aogami sẽ mãi mãi nằm trong bóng tối.

Mọi thứ yên bình đến mức tôi gần như đinh ninh rằng nhiệm vụ của chúng tôi tại thế giới này đã hoàn thành, và đây là lúc chúng tôi có thể rời khỏi thị trấn này để quay về thế giới thực. Đêm nay, khi đứng ngoài ban công nhìn xuống thị trấn im lìm dưới ánh trăng, tôi khẽ chạm vào mặt đồng hồ.

``\eng{Này Game-san, ông đã tìm được cổng ra hay bất kỳ lỗ xoáy không gian nào chưa?}'' tôi hỏi với giọng nhỏ nhẹ.

Một khoảng lặng ngắn vang lên trước khi giọng Game đáp lại, mang theo sự trầm tư hiếm thấy: ``\eng{Chưa, Kuon-user. Tôi đã quét toàn bộ khu vực thị trấn và rừng Aogami trong bán kính 20km nhưng không phát hiện bất kỳ dao động nào của cánh cổng dịch chuyển.}''

``Sao?'' tôi cau mày, cảm giác bất an len lỏi trong lòng. ``Kỳ lạ thật. Nghi thức siêu thoát đã hoàn hảo, oán khí đã tan biến\dots\ Theo logic thông thường, cánh cửa để chúng ta quay về thế giới thực phải xuất hiện rồi chứ?''

Game không đáp lại ngay. Sự im lặng của ông ấy khiến tôi nhận ra rằng, có lẽ, câu chuyện ở Aogami vẫn chưa thực sự khép lại như tôi tưởng. Chúng tôi\dots\ vẫn còn một bước nào đó vẫn chưa hoàn thành sao? Hay thị trấn này vẫn còn một điều gì đó chưa được giải quyết?

\pov{-}

Đêm khuya, khi ánh đèn trong nhà Hikawa đã tắt lịm và tiếng thở đều đặn của ba anh em chìm sâu vào giấc ngủ. Khu vườn phía sau nhà tĩnh lặng tuyệt đối.

Dưới ánh trăng bạc, bóng của những tán cây in dài trên mặt đất. Một bóng người nhỏ nhắn bất chợt xuất hiện bên ngoài cửa sổ phòng khách. Cô bé đứng đó, lặng lẽ nhìn vào bên trong.

Mái tóc đen tết đuôi sam khẽ đung đưa theo làn gió đêm. Đôi mắt xanh lục trong veo như ngọc bích, phản chiếu ánh trăng lấp lánh. Nhưng khác với linh hồn mờ ảo từng xuất hiện ở rừng Aogami, cô bé lúc này mang một vẻ thực tại đến lạ lùng. Dưới chân em, một cái bóng đổ dài trên thảm cỏ, rõ nét và vững chãi. Một thực thể bằng xương bằng thịt.

Cô bé khẽ nghiêng đầu, đôi môi nhỏ nhắn mấp máy tạo thành một nụ cười dịu dàng.

``Lần này\dots'' giọng nói trong trẻo như tiếng chuông ngân vang lên giữa màn đêm tĩnh lặng, ``\dots\ em sẽ không biến mất nữa đâu. Chúng ta sẽ sớm gặp lại nhau thôi, những người bạn của em.''

Cô bé lùi lại một bước, biến mất vào bóng tối của khu vườn, để lại một lời hứa còn vương vất trong gió, báo hiệu cho một hành trình mới sắp bắt đầu.

\end{document}